\BourbakiExercisesSection

\begin{enumerate}
\def\labelenumi{\arabic{enumi}.}
\item
  枚举长度 \(\leqslant 4\) 的、位于 \(\mathbf{M}(\mathbf{X})\) 中的元素（此时 \(\mathbf{X}\) 由单个元素组成）。
\item
  设 X 为一个集合，M 为广群 M(X)、Mo(X) 或 \(\mathrm{N}^{(X)}\) 之一。证明 M 的每个自同构都使 X 稳定。由此推出一个从 \(\mathfrak{S}_{X}\) 到群 Aut(M) 上的同构。
\end{enumerate}

证明 M 的自同态与从 X 到 M 的映射一一对应。

\begin{enumerate}
\def\labelenumi{\arabic{enumi}.}
\setcounter{enumi}{2}
\tightlist
\item
  设 X 为一个集合， \(M_{a}\) 为在单元素集 \(\{a\}\) 上构造的自由广群。设 \(\rho\) 表示将 X 的每个元素都映到 a 的、从 \(\mathrm{M}(\mathrm{X})\) 到 \(M_{a}\) 中的同态；另一方面，设 \(\lambda: \mathrm{M}(\mathrm{X}) \to \mathrm{Mo}(\mathrm{X})\) 为第 9 号中定义的同态。证明映射 \(w \mapsto (\lambda(w), \rho(w))\) 是从 \(\mathrm{M}(\mathrm{X})\) 到 \(\mathrm{Mo}(\mathrm{X}) \times \mathrm{M}_{a}\) 的由 u 与 v 具有相同长度的有序对 \((u, v)\) 所组成的子广群上的同构。
\end{enumerate}

¶ 4. 设 X 为一个集合，c 为不属于 X 的元素。记 Y = X ∪ \{c\}。设 M 为以 Mo(Y) 为底层集合、合成律由 (u, v) ↦ cuv 给出的广群。设 L 为由 X 生成的 M 的子广群，并设 ε 为从 Y 到 Z 的映射，它在 X 上取值 -1，在 c 处取值 1。

\begin{enumerate}
\def\labelenumi{(\alph{enumi})}
\tightlist
\item
  证明 L 是由满足下列两个条件的字 \(w = y_{1} \ldots y_{p}\) , \(y_{i} \in Y\) 组成的 M 的子集：
\end{enumerate}

\[
(i) \varepsilon (y _ {1}) + \dots + \varepsilon (y _ {p}) = - 1
\]

\[
(ii) \varepsilon \left(y _ {1}\right) + \dots + \varepsilon \left(y _ {i}\right) \geqslant 0 \quad \text { for } 1 \leqslant i \leqslant p - 1.
\]

证明 L - X 的每个元素都可以唯一地写成 cuv 的形式，其中 \(u, v \in L\) 。

（满足 (i) 与 (ii) 的元素 \(w\) 构成一个子广群 \(L'\) ，它含于 \(M\) 且包含 \(X\) 。若 \(w = y_1 \ldots y_p\) 属于 \(L'\) ，设 \(k\) 为使 \(\varepsilon(y_1) + \cdots + \varepsilon(y_k) = 0\) 的最小整数；证明 \(y_1 = c\) ，且元素 \(u = y_2 \ldots y_k\) 与 \(v = y_{k+1} \ldots y_p\) 属于 \(L'\) ；于是 \(w = cuv\) ，从而对 \(p\) 作归纳即得关系 \(w \in L\) 。然后证明，关系 \(u', v' \in L\) 与 \(w = cu'v'\) 蕴含 \(u' = u, v' = v\) 。）

\begin{enumerate}
\def\labelenumi{(\alph{enumi})}
\setcounter{enumi}{1}
\tightlist
\item
  证明 X 到 L 中的单射可延拓为 M(X) 到 L 上的同构。
\end{enumerate}

\begin{enumerate}
\def\labelenumi{\arabic{enumi}.}
\setcounter{enumi}{4}
\tightlist
\item
  设 X 为单元素集。若 n 为整数 \(\geqslant1\) ，设 \(u_{n}\) 表示 M(X) 中长度为 n 的元素个数。
\end{enumerate}

\begin{enumerate}
\def\labelenumi{(\alph{enumi})}
\tightlist
\item
  对于每一个集合Y，证明
\end{enumerate}

\[
\operatorname{Card} \left(\mathrm{M} _ {n} (\mathrm{Y})\right) = u _ {n}. \operatorname{Card} (\mathrm{Y}) ^ {n}.
\]

（利用习题 3。）

\begin{enumerate}
\def\labelenumi{(\alph{enumi})}
\setcounter{enumi}{1}
\tightlist
\item
  建立关系
\end{enumerate}

\[
u _ {n} = \sum_ {p = 1} ^ {n - 1} u _ {p} u _ {n - p} \quad \text { for } n \geqslant 2.
\]

\begin{enumerate}
\def\labelenumi{(\alph{enumi})}
\setcounter{enumi}{2}
\tightlist
\item
  *设 \(f(\mathbf{T})\) 是形式幂级数 \(\sum_{n=1}^{\infty} u_n \mathbf{T}^n\) 。证明
\end{enumerate}

\[
f (\mathbf {T}) = \mathbf {T} + f (\mathbf {T}) ^ {2}.
\]

由此推出公式

\[
f (\mathrm{T}) = \frac {1}{2} (1 - (1 - 4 \mathrm{T}) ^ {1 / 2}), \quad u _ {n} = \frac {2 ^ {n - 1}}{n !} 1. 3. 5 \dots (2 n - 3) \quad \text { for } n \geqslant 2. *
\]

利用《集合论》III，§ 5，习题 11 得到最后这一结果。

¶ 6. 设 X 为一个集合。

\begin{enumerate}
\def\labelenumi{(\alph{enumi})}
\item
  设 N 为 M(X) 的子广群，且 Y = N - (N.N)。单射 Y → N 可延拓为一个同态 u: M(Y) → N。证明 u 是同构。（换言之，自由广群的每个子广群都是自由的。）
\item
  若 \(x \in \mathbf{M}(\mathbf{X})\) ，设 \(\mathbf{M}_x\) 表示 \(\mathbf{M}(\mathbf{X})\) 中由 \(x\) 生成的子广群；由 (a)， \(\mathbf{M}_x\) 等同于在 \(x\) 上构造的自由广群。证明，若 \(x, y \in \mathbf{M}(\mathbf{X})\) ，则要么 \(\mathbf{M}_x \subset \mathbf{M}_y\) ，要么 \(\mathbf{M}_y \subset \mathbf{M}_x\) ，要么 \(\mathbf{M}_x \cap \mathbf{M}_y = \varnothing\) 。（若 \(\mathbf{M}_x \cap \mathbf{M}_y \neq \varnothing\) ，设 \(z\) 为 \(\mathbf{M}_x \cap \mathbf{M}_y\) 中长度最小的元素；证明 \(z = x\) 或 \(z = y\) 。）
\end{enumerate}

\begin{enumerate}
\def\labelenumi{\arabic{enumi}.}
\setcounter{enumi}{6}
\tightlist
\item
  设 \(\mathbf{M}_x\) 为在单元素集 \(\{x\}\) 上构造的自由广群。(a) 对所有 \(y \in \mathbf{M}_x\) ，证明存在唯一的自同态 \(f_y\) （\(\mathbf{M}_x\) 的），使得 \(f_y(x) = y\) ；它是从 \(\mathbf{M}_x\) 到子广群 \(\mathbf{M}_y\) （由 \(y\) 生成）上的同构（参见习题 6）。
\end{enumerate}

\begin{enumerate}
\def\labelenumi{(\alph{enumi})}
\setcounter{enumi}{1}
\item
  若 \(y, z \in \mathbf{M}_x\) ，记 \(y \circ z = f_y(z)\) 。证明合成律 \((y, z) \mapsto y \circ z\) 使 \(\mathbf{M}_x\) 成为一个以 \(x\) 为单位元的幺半群，它同构于幺半群 \(\operatorname{End}(\mathbf{M}_x)\) ——\(\mathbf{M}_x\) 的自同态所组成的幺半群。
\item
  设 \(\mathbf{N}\) 为 \(\mathbf{M}_x\) 的异于 \(\mathbf{M}_x\) 的单生成子广群组成的集合。一个元素 \(y\) （属于 \(\mathbf{M}_x\) ）称为本原的，如果 \(\mathbf{M}_y\) 是 \(\mathbf{N}\) 的一个极大元。证明 \(xx, x(x(xx))\) 是本原的，而 \(x, (xx)(xx)\) 不是。
\item
  设 \(\mathbf{P}\) 为 \(\mathbf{M}_x\) 的本原元组成的集合。证明，若 \(y, z \in \mathbf{P}, y \neq z\) ，则 \(\mathbf{M}_y \cap \mathbf{M}_z = \varnothing\) （利用习题 6）。
\item
  设 \(\operatorname{Mo}(\mathbf{P})\) 为在 \(\mathbf{P}\) 上构造的自由幺半群。若给 \(\mathbf{M}_x\) 赋予 (b) 中定义的幺半群结构，则单射 \(\mathbf{P} \to \mathbf{M}_x\) 可延拓为一个同态 \(p: \operatorname{Mo}(\mathbf{P}) \to \mathbf{M}_x\) 。证明 \(p\) 是一个同构。
\end{enumerate}

（若 \(z \in \mathbf{M}_x\) ，通过对 \(l(z)\) 作归纳证明 \(z \in \operatorname{Im}(p)\) 。另一方面，若 \(y_1, \ldots, y_n, z_1, \ldots, z_m \in \mathbb{P}\) 满足 \(p(y_1 \ldots y_n) = p(z_1 \ldots z_m)\) ，则 \(\mathbf{M}_{y_1} \cap \mathbf{M}_{z_1} \neq \varnothing\) ，从而 \(y_1 = z_1\) 且 \(f_{y_1}(p(y_2 \ldots y_n)) = f_{y_1}(p(z_2 \ldots z_m))\) ；由此得 \(p(y_2 \ldots y_n) = p(z_2 \ldots z_m)\) ，从而推出 \(y_1 \ldots y_n = z_1 \ldots z_m\) （对 \(\sup(n, m)\) 作归纳）。）

\begin{enumerate}
\def\labelenumi{\arabic{enumi}.}
\setcounter{enumi}{7}
\tightlist
\item
  设 X 为一个集合；对每个整数 \(q \geqslant 0\) ，设 \(\mathrm{Mo}(\mathbf{X})_{q}\) 为 \(\mathrm{Mo}(\mathbf{X})\) 中长度为 q 的元素组成的集合，且设 \(\mathrm{Mo}^{(q)}(\mathbf{X})\) 为诸 \(\mathrm{Mo}(\mathbf{X})_{nq}, n \in \mathbb{N}\) 的并。单射 \(\mathrm{Mo}(\mathbf{X})_{q} \to \mathrm{Mo}^{(q)}(\mathbf{X})\) 可延拓为一个同态
\end{enumerate}

\[
\operatorname{Mo} (\operatorname{Mo} (\mathrm{X}) _ {q}) \rightarrow \operatorname{Mo} ^ {(q)} (\mathrm{X});
\]

证明当 \(q \geqslant 1\) 时该同态是双射。

\begin{enumerate}
\def\labelenumi{\arabic{enumi}.}
\setcounter{enumi}{8}
\item
  设 \(x, y \in \mathbf{X}\) 且 \(x \neq y\) 。证明 \(\operatorname{Mo}(\mathbf{X})\) 中由 \(\{x, xy, yx\}\) 生成的子幺半群不同构于任何自由幺半群。
\item
  证明由表现
\end{enumerate}

\[
\langle x, y; x y ^ {2} = y ^ {3} x, y x ^ {2} = x ^ {3} y \rangle
\]

定义的群约化为单位元。

（第一个关系蕴含 \(x^{2}y^{8}x^{-2} = y^{18}\) 与 \(x^{3}y^{8}x^{-3} = y^{27}\) ；利用第二个关系可推出 \(y^{18} = y^{27}\) ，从而 \(y^{9} = e\) ；于是 \(y^{2}\) 共轭于 \(y^{3}\) 这一事实蕴含 \(y = e\) ，从而 \(x = e\) 。）

\begin{enumerate}
\def\labelenumi{\arabic{enumi}.}
\setcounter{enumi}{10}
\tightlist
\item
  由表现 \(\langle x,y;x^{2},y^{3},(xy)^{2}\rangle\) 定义的群同构于 \(S_{3}\) 。
\end{enumerate}

¶ 12. 由表现 \(\langle x, y; x^3, y^2, (xy)^3 \rangle\) 定义的群同构于 \(\mathfrak{A}_4\) 。（证明 \(y\) 的共轭元与单位元一起构成一个阶 \(\leqslant 4\) 、指数 \(\leqslant 3\) 的正规子群。）

\begin{enumerate}
\def\labelenumi{\arabic{enumi}.}
\setcounter{enumi}{12}
\tightlist
\item
  设 \(G\) 为一个群，\(X\) 为 \(G\) 的一个与 \(X^{-1}\) 不相交的子集。设 \(S = X \cup X^{-1}\) 。证明下列两个性质的等价性：
\end{enumerate}

\begin{enumerate}
\def\labelenumi{(\roman{enumi})}
\item
  X 是 G 中的自由族。
\item
  没有乘积 \(s_1 \ldots s_n\) （满足 \(n \geqslant 1\) , \(s_i \in S\) 且 \(s_i s_{i+1} \neq e\) 对 \(1 \leqslant i \leqslant n-1\) 成立）等于 \(e\) 。
\end{enumerate}

\begin{enumerate}
\def\labelenumi{\arabic{enumi}.}
\setcounter{enumi}{13}
\tightlist
\item
  群 G 称为自由群，如果它拥有一个基本族，从而同构于在某个集合 X 上构造的自由群 F(X)。
\end{enumerate}

\begin{enumerate}
\def\labelenumi{(\alph{enumi})}
\tightlist
\item
  证明，若 \((t_i)_{i\in I}\) 与 \((t_j')_{j\in J}\) 是两个这样的族，则 \(\operatorname{Card}(I) = \operatorname{Card}(J)\) 。（当 \(\operatorname{Card}(I)\) 无限时，证明 \(\operatorname{Card}(J)\) 也无限，且二者都等于 \(\operatorname{Card}(G)\) 。当 \(\operatorname{Card}(I)\) 是整数 \(d\) 时，证明 \(G\) 中指数为 2 的子群个数是 \(2^d - 1\) 。）
\end{enumerate}

G 的基本族的基数称为 G 的秩。

\begin{enumerate}
\def\labelenumi{(\alph{enumi})}
\setcounter{enumi}{1}
\item
  一个自由群的秩为 0（相应地，1）当且仅当它约化为 \(\{e\}\) （相应地，同构于 \(\mathbf{Z}\) ）。
\item
  设 G 为秩为 d 的自由群，且 \((x_{1},\ldots,x_{d})\) 为 G 的由 d 个元素组成的生成族。证明该族是基本的。（利用习题 34 与 § 5 习题 5。）
\item
  证明秩 \(\geqslant 2\) 的自由群包含一个具有给定秩 \(d\) 的自由子群，这里 \(d \leqslant \aleph_0\) 为任意基数。（利用习题 22。）
\end{enumerate}

¶ 15. 设 G 是一个具有表现 (t, r) 的群，其中 \(\mathbf{t} = (t_{i})_{i \in I}\) 与 \(\mathbf{r} = (r_{j})_{j \in J}\) ；令 \(\mathrm{F}((\mathrm{T}_{i})_{i \in I})\) 表示自由群 \(\mathrm{F}(I)\)，并令 \(\mathbf{T} = (\mathrm{T}_{i})\) 。另一方面，设 \(x = (x_{\alpha})_{\alpha \in A}\) 是 G 的一个生成族；记

\[
\mathrm{F} (\mathbf {A}) = \mathrm{F} ((\mathbf {X} _ {\alpha}) _ {\alpha \in A}),
\]

\(\mathbf{X} = (\mathbf{X}_{\alpha})\) ，而 \(\mathbf{R}_{\mathbf{x}}\) 表示 \(\mathbf{x}\) 的关系子集；它是 \(\mathrm{F(A)}\) 的一个正规子群。对所有 \(i \in \mathbf{I}\) ，设 \(\phi_i(\mathbf{X})\) 为 \(\mathrm{F(A)}\) 中使 \(t_i = \phi_i(\mathbf{x})\) 在 \(\mathbf{G}\) 中成立的元素；对所有 \(\alpha \in \mathbf{A}\) ，设 \(\psi_{\alpha}(\mathbf{T})\) 为 \(\mathrm{F(I)}\) 中使 \(x_{\alpha} = \psi_{\alpha}(\mathbf{t})\) 在 \(\mathbf{G}\) 中成立的元素。证明 \(\mathbf{R}_{\mathbf{x}}\) 是 \(\mathrm{F(A)}\) 中由元素 \(\mathbf{X}_{\alpha}^{-1} \psi_{\alpha}(\phi_i(\mathbf{X})_{i \in I}), \alpha \in \mathbf{A}\) 与 \(r_j(\phi_i(\mathbf{X})_{i \in I}) j \in J\) 生成的正规子群。（若 \(\mathbf{R}_{\mathbf{x}}'\) 表示 \(\mathrm{F(A)}\) 中由上述元素生成的正规子群，则 \(\mathbf{R}_{\mathbf{x}}' \subset \mathbf{R}_{\mathbf{x}}\) ；另一方面，证明同态 \(\mathrm{F(I)} \to \mathrm{F(A)}\) （由诸 \(\phi_i\) 所定义）在过渡到商时给出一个同态 \(\mathrm{G} \to \mathrm{F(A)} / \mathrm{R}_{\mathbf{x}}'\) ，它是典范同态 \(\mathrm{F(A)} / \mathrm{R}_{\mathbf{x}}' \to \mathrm{F(A)} / \mathrm{R}_{\mathbf{x}} = \mathrm{G}\) 的逆。）

\begin{enumerate}
\def\labelenumi{\arabic{enumi}.}
\setcounter{enumi}{15}
\tightlist
\item
  一个群称为有限生成的（相应地，有限表现的），如果它有一个表现 \(\langle (t_i)_{i\in I}, (r_j)_{j\in J}\rangle\) ，其中 \(I\) 是有限集（相应地，其中 \(I\) 与 \(J\) 都是有限集）。
\end{enumerate}

\begin{enumerate}
\def\labelenumi{(\alph{enumi})}
\item
  设 G 是一个有限表现的群，且 \(\mathbf{x} = (x_{\alpha})_{\alpha \in \mathbf{A}}\) 是 G 的一个有限生成族。证明存在族 x 的关系子的一个有限族 \(\mathbf{s} = (s_{k})_{k \in \mathbf{K}}\) ，使得 \((\mathbf{x}, \mathbf{s})\) 是 G 的一个表现。（利用前一习题。）
\item
  证明有限群是有限表现的。
\item
  给出一个有限表现的群的例子，它具有一个不是有限生成的子群。
\item
  若 \(\mathbf{G}\) 是一个有限表现的群，而 \(\mathbf{H}\) 是 \(\mathbf{G}\) 的一个正规子群，证明下列性质的等价性：
\item
  G/H 是有限表现的。
\end{enumerate}

\begin{enumerate}
\def\labelenumi{(\roman{enumi})}
\setcounter{enumi}{1}
\item
  存在 G 的有限子集 X，使得 H 是 G 中由 X 生成的正规子群。
\item
  若 \((\mathbf{H}_i)_{i\in \mathbf{I}}\) 是 \(\mathbf{G}\) 的一个右有向正规子群族，其并等于 \(\mathbf{H}\) ，则存在 \(i\in \mathbf{I}\) ，使得 \(\mathbf{H}_i = \mathbf{H}\) 。
\end{enumerate}

（利用 (a) 证明 (i) 蕴含 (ii)。）

\begin{enumerate}
\def\labelenumi{(\alph{enumi})}
\setcounter{enumi}{4}
\item
  若 G 是一个有限表现的群，证明 \(\mathbf{G}/\mathbf{C}^{r}(\mathbf{G})\) 对所有 \(r \geqslant 1\) 都是有限表现的。（利用 § 6 习题 15。）
\item
  设 \((\mathbf{G}_{\alpha})_{\alpha \in \mathbf{A}}\) 为一族群，且设 \(\mathbf{G} = \prod_{\alpha \in \mathbf{A}} \mathbf{G}_{\alpha}\) 为诸 \(\mathbf{G}_{\alpha}\) 的积。证明 \(\mathbf{G}\) 是有限生成的（相应地，有限表现的）当且仅当每个 \(\mathbf{G}_{\alpha}\) 都是有限生成的（相应地，有限表现的），且 \(\mathbf{G}_{\alpha} = \{e\}\) 对除有限个以外的 \(\alpha\) 成立。
\end{enumerate}

¶ 17. (a) 证明有限生成的幂零群的每一个子群都是有限生成的。（在交换情形中，通过对群的最小生成元个数进行归纳来论证；然后对群的幂零类进行归纳。）

给出一个有限生成的可解群的例子，它包含一个不是有限生成的子群。

\begin{enumerate}
\def\labelenumi{(\alph{enumi})}
\setcounter{enumi}{1}
\tightlist
\item
  证明每个有限生成的幂零群都是有限表现的。（将该群写成形式 \(\mathrm{F}(\mathbf{X}) / \mathbb{R}\) ，其中 \(\mathbf{X}\) 有限，并选取 \(r\) 使得 \(\mathbf{R} \supset \mathbf{C}^r (\mathrm{F}(\mathbf{X}))\) ；利用 (a) 证明 \(\mathbb{R} / \mathbb{C}^r (\mathrm{F}(\mathbf{X}))\) 是有限生成的；注意 \(\mathrm{F}(\mathbf{X}) / \mathbb{C}^r (\mathrm{F}(\mathbf{X}))\) 是有限表现的，参见习题 16。）
\end{enumerate}

\begin{enumerate}
\def\labelenumi{\arabic{enumi}.}
\setcounter{enumi}{17}
\tightlist
\item
  设 S 是群 G 的一个对称子集。G 的两个元素 \(g_{1}, g_{2}\) 称为 S-邻元，如果 \(g_{2}^{-1}g_{1}\) 属于 S。G 的元素的一个序列 \((g_{1}, \ldots, g_{n})\) 称为 S-链，如果 \(g_{i}\) 与 \(g_{i+1}\) 是 S-邻元（对 \(1 \leqslant i \leqslant n - 1\) ）；元素 \(g_{1}\) 与 \(g_{n}\) 分别称为链的起点与终点。
\end{enumerate}

\begin{enumerate}
\def\labelenumi{(\alph{enumi})}
\item
  设 Y 是 G 的一个子集，\(R_{Y}\{a, b\}\) 表示关系“在 Y 中存在一条从 a 开始而在 \(b''\) 结束的 S-链”。证明 \(R_{Y}\) 是 Y 上的一个等价关系（参见《集合论》II，§ 6，习题 10）；该关系下的等价类称为 Y 的连通 S-分量。若 Y 至多有一个连通分量，则称 Y 是 S-连通的。
\item
  证明 \(\mathbf{G}\) 是 S-连通的当且仅当 \(\mathbf{S}\) 生成 \(\mathbf{G}\) 。
\item
  假设没有元素 \(s \in S\) 满足关系 \(s^2 = e\) ；选取子集 \(X\) （含于 \(S\) ），使得 \(S\) 是 \(X\) 与 \(X^{-1}\) 的不交并。证明下列条件的等价性：
\item
  族 \(\mathbf{X}\) 是自由的。
\end{enumerate}

\begin{enumerate}
\def\labelenumi{(\roman{enumi})}
\setcounter{enumi}{1}
\tightlist
\item
  不存在 G 中的 S-链 \((g_{1},\ldots ,g_{n})\) ——由 \(n\geqslant 3\) 个互异元素组成——使得 \(g_{1}\) 与 \(g_{n}\) 是 S-邻元。
\end{enumerate}

（利用习题 13。）

¶ 19. 设 X 为一个集合，G = F(X) 为在 X 上构造的自由群，且 S = X ∪ X \(^{-1}\) 。

\begin{enumerate}
\def\labelenumi{(\alph{enumi})}
\item
  若 \(g \in G\) ，每个序列 \((s_1, \ldots, s_n)\) （由 \(S\) 的元素组成，满足 \(g = s_1 \ldots s_n\) 且 \(s_i s_{i+1} \neq e\) 对 \(1 \leqslant i < n\) 成立）称为 \(g\) 的既约分解。证明 \(G\) 的每个元素有且仅有一个既约分解；整数 \(n\) 称为 \(g\) 的长度（参见习题 26）。
\item
  设 Y 是 G 的包含 e 的子集。证明以下条件的等价性：
\end{enumerate}

\((b_{1})\) Y 是 S-连通的（参见习题 18）。

\((b_{2})\) 对所有 \(g \in \mathbf{Y}\) ，若 \((s_1, \ldots, s_n)\) 是 \(g\) 的既约分解，则 \(s_1 \ldots s_i \in \mathbf{Y}\) 对 \(1 \leqslant i \leqslant n\) 成立。

\begin{enumerate}
\def\labelenumi{(\alph{enumi})}
\setcounter{enumi}{2}
\tightlist
\item
  设 \(Y_{1}\) 与 \(Y_{2}\) 是 G 的两个非空的 S-连通子集，且 \(Y_{1} \cap Y_{2} = \varnothing\) 。证明至多存在一个有序对
\end{enumerate}

\[
\left(y _ {1}, y _ {2}\right) \in \mathbf {Y} _ {1} \times \mathbf {Y} _ {2}
\]

使得 \(y_{1}\) 与 \(y_{2}\) 是 S-邻元。这样的有序对存在当且仅当 \(\mathbf{Y}_1 \cup \mathbf{Y}_2\) 是 S-连通的；这时称 \(\mathbf{Y}_1\) 与 \(\mathbf{Y}_2\) 相邻。

\begin{enumerate}
\def\labelenumi{(\alph{enumi})}
\setcounter{enumi}{3}
\tightlist
\item
  设 \(Y_{1}, \ldots, Y_{n}\) 是 G 的一列互不相交的非空 S-连通子集。假设 \(Y_{i}\) 与 \(Y_{i+1}\) 相邻（对 \(1 \leqslant i < n\) ）。证明，若 \(n \geqslant 3\) ，则 \(Y_{1}\) 与 \(Y_{n}\) 不相邻。
\end{enumerate}

¶ 20. 我们保留前一习题的记号。设 H 是 G 的一个子群。

\begin{enumerate}
\def\labelenumi{(\alph{enumi})}
\item
  设 \(R_{H}\) 为 G 中满足 \(hT \cap T = \varnothing\) （若 \(h \in H, h \neq e\) ）的 S-连通子集 T 组成的集合。证明 \(R_{H}\) 按包含关系是归纳的。
\item
  设 Y 为 \(R_{H}\) 的一个极大元。证明 \(G = \bigcup_{h \in H} hY\) ，换言之，Y 是模 H 的右陪集的一个代表系。（若 \(G' = \bigcup hY\) 不同于 G，证明存在 \(x \in G'\) 与 \(y \in Y\) ，二者是 S-邻元，并由此推出 \(Y \cup \{x\}\) 属于 \(R_{H}\) 。）
\item
  设 \(S_{H}\) 为由使 Y 与 hY 相邻的元素 \(h \in H\) 组成的集合。它是 H 的一个对称子集。证明元素 \(h \in H\) 属于 \(S_{H}\) 当且仅当 \(h \neq e\) 且存在 \(y, y' \in Y, s \in S\) 使得 \(ys = hy'\) 。
\end{enumerate}

证明 \(S_{H}\) 生成 H。（若 \(H'\) 是由 \(S_{H}\)

生成的 H 的子群，证明 \(\bigcup_{h\in H'}hY\) 是 G 的一个连通 S-分量。）通过 § 4 习题 14 得到这一结果。

\begin{enumerate}
\def\labelenumi{(\alph{enumi})}
\setcounter{enumi}{3}
\tightlist
\item
  设 \(X_{H}\) 为 \(S_{H}\) 的一个子集，使得 \(S_{H}\) 是 \(X_{H}\) 与 \(X_{H}^{-1}\) 的不交并（证明这样的子集存在）。
\end{enumerate}

证明 \(X_{H}\) 是 H 的一个基本族，特别地，H 是自由群（“尼尔森–施赖尔定理”）。（将习题 18(c) 的判据应用于 H 及 \(S_{H}\) ，注意 H 的两个元素 h 与 \(h'\) 是 \(S_{H}\) -邻元当且仅当 hY 与 \(h'Y\) 是 G 的相邻子集。利用习题 19(d) 证明 \(S_{H}\) 满足习题 18(c) 中的条件 (ii)。）

\begin{enumerate}
\def\labelenumi{(\alph{enumi})}
\setcounter{enumi}{4}
\tightlist
\item
  设 \(d = (\mathbf{G}:\mathbf{H}) = \operatorname{Card}(\mathbf{Y})\) 。假设 \(d\) 是有限的。证明：
\end{enumerate}

\[
\begin{array}{l l} \operatorname{Card} (\mathbf {X} _ {\mathrm{H}}) = \operatorname{Card} (\mathbf {X}) = \operatorname{Card} (\mathbf {G}) & \text { if } \operatorname{Card} (\mathbf {X}) \text { is   infinite } \\ \operatorname{Card} (\mathbf {X} _ {\mathrm{H}}) = 1 + d (\operatorname{Card} (\mathbf {X}) - 1) & \text { if } \operatorname{Card} (\mathbf {X}) \text { is   finite }. \end{array}
\]

（设 T 是 G 的一个非空 S-连通有限子集，T'（相应地 T''）是其第一个元素（相应地两个元素）属于（相应地都属于）T 的 S-相邻元素的有序对组成的集合。Card(T') = 2Card(X)Card(T)，且应通过对 Card(T) 作归纳证明

\[
\operatorname{Card} \left(\mathrm{T} ^ {\prime \prime}\right) = 2 \operatorname{Card} (\mathrm{T}) - 2.
\]

将结果应用于 T = Y，注意 \(S_{H}\) 与 \(T' - T''\) 等势。

\begin{enumerate}
\def\labelenumi{\arabic{enumi}.}
\setcounter{enumi}{20}
\item
  设 H 是有限指数的群 G 的一个子群。证明 H 是有限表现的当且仅当 G 是有限表现的。（可假定 G 是有限生成的，参见 § 4 习题 14。将 G 写成形式 G = F/R，其中 F 是自由且有限生成的；H 在 F 中的原像 F' 是自由且有限生成的，参见习题 20，且 H = F'/R。若 R 作为 F 的正规子群由 \((r_{j})_{j\in J}\) 生成，证明它作为 F' 的正规子群由诸 \(yr_{j}y^{-1}\) , \(y\in Y\) , \(j\in J\) 生成，其中 Y 表示 F 中模 F' 的右陪集的一个代表系。）
\item
  \begin{enumerate}
  \def\labelenumii{(\alph{enumii})}
  \tightlist
  \item
    设 \((y_{n})_{n\in \mathbf{Z}}\) 是群 F 的一个基本族。设 \(\phi\) 为将 \(y_{n}\) 映为 \(y_{n + 1}\) 的 F 的自同构；该自同构定义了运算 \(\tau\) （\(\mathbf{Z}\) 在 F 上的运算）；设 \(\mathbf{E} = \mathbf{F}\times_{\tau}\mathbf{Z}\) 为相应的半直积。证明元素 \(x = (e,1)\) 与 \(y = (y_0,0)\) 构成 E 的一个基本族。（设 \(\mathrm{F}_{x,y}\) 为在 \(\{x,y\}\) 上构造的自由群，并设 \(f\) 为 \(\mathrm{F}_{x,y}\) 到 E 中的典范同态。证明存在同态 \(g\colon \mathbf{E}\to \mathbf{F}_{x,y}\) ，使得 \(g((0,1)) = x\) 且 \(g((y_n,0)) = x^n yx^{-n}\) ，并且 \(f\) 与 \(g\) 互为逆映射。）
  \end{enumerate}
\end{enumerate}

\begin{enumerate}
\def\labelenumi{(\alph{enumi})}
\setcounter{enumi}{1}
\item
  由此推出，\(F_{x,y}\) 的由 y 生成的正规子群是一个自由群，其基本族为 \((x^{n}yx^{-n})_{n\in\mathbf{Z}}\) 。将习题 20 应用于由诸 \(x^{n}, n\in Z\) 组成的代表系来得到这一结果。
\item
  将上述结论推广到任意秩的自由群。
\end{enumerate}

\begin{enumerate}
\def\labelenumi{\arabic{enumi}.}
\setcounter{enumi}{22}
\tightlist
\item
  设 \(\mathbf{F}\) 是基本族为 \((x, y)\) , \(x \neq y\) 的一个自由群。证明 F 的导群以换位子族 \((x^i, y^j)\) , \(i \in \mathbf{Z}\) , \(j \in \mathbf{Z}\) , \(i \neq 0\) , \(j \neq 0\) 为基本族。
\end{enumerate}

（利用习题 20 或习题 32。）

\begin{enumerate}
\def\labelenumi{\arabic{enumi}.}
\setcounter{enumi}{23}
\tightlist
\item
  设 G 是一个群，\(S = G - \{e\}\) 且 \(F = F(S)\) 。对所有 \(s \in S\) ，设 \(x_s\) 表示 F 中相应的元素。若 \(s, t \in S\) ，设 \(r_{s,t}\) 表示由下列方式定义的 F 中的元素：
\end{enumerate}

\[
r _ {s, t} = x _ {s} x _ {t} x _ {s t} ^ {- 1} \quad \text { if } s t \neq e \text { in } G
\]

\[
r _ {s, t} = x _ {s} x _ {t} \quad \text { if } s t = e \text { in } G.
\]

设 \(\phi\) 为将 \(x_{s}\) 映为 \(s\) 的、F 到 G 上的同态，并令 R 为 \(\phi\)

的核。证明族 \((r_{s,t})\) （\(s, t \in \mathbf{S} \times \mathbf{S}\)）是群 \(\mathbf{R}\) 的一个基本族。（将习题 20 应用于 \(\mathbf{G}\) 在 \(\mathbf{F}\) 中的由 \(e\) 与诸 \(x_{s'}\) 组成的代表系。）

\begin{enumerate}
\def\labelenumi{\arabic{enumi}.}
\setcounter{enumi}{24}
\tightlist
\item
  设 G 是一个群。证明以下性质的等价性：
\end{enumerate}

\begin{enumerate}
\def\labelenumi{(\alph{enumi})}
\item
  G 是一个自由群。
\item
  对于每个群 H 以及 G 被 H 的扩张 E，存在一个截面 \(G \rightarrow E\) .
\end{enumerate}

（为证明 (b) 蕴含 (a)，取 E 为一个自由群并使用习题 20。）

\begin{enumerate}
\def\labelenumi{\arabic{enumi}.}
\setcounter{enumi}{25}
\tightlist
\item
  设 \(\mathbf{X}\) 为一个集合，\(g\) 为自由群 \(\mathrm{F}(\mathbf{X})\) 的一个元素。设 \(g = \prod_{\alpha=1}^{n} x_{\alpha}^{m(\alpha)}\) 为 \(g\) 作为 \(\mathbf{X}\) 中元素的幂之乘积的典范分解（参见第 5 号，命题 7）。整数 \(l(g) = \sum_{\alpha} |m(\alpha)|\) 称为 \(g\) 的长度。
\end{enumerate}

\begin{enumerate}
\def\labelenumi{(\alph{enumi})}
\item
  \(g\) 称为循环既约的，如果或者 \(x_{1} \neq x_{n}\) ，或者 \(x_{1} = x_{n}\) 且 \(m(1)m(n) > 0\) 。证明 \(F(X)\) 的每个共轭类包含且仅包含一个循环既约元素；它是所讨论的类中长度最小的元素。
\item
  一个元素 \(x \in \mathrm{F}(\mathbf{X})\) 称为本原的，如果不存在整数 \(n \geqslant 2\) 和元素 \(g \in \mathrm{F}(\mathbf{X})\) 使得 \(x = g^n\) 。证明每个元素 \(z \neq e\) 在 \(\mathrm{F}(\mathbf{X})\) 中都可以唯一地写成形式 \(x^n\) ，其中 \(n \geqslant 1\) 且 \(x\) 是本原的。（将其归约为 \(z\) 是循环既约的情形，并注意到，若 \(z = x^n\) ，则元素 \(x\) 也是循环既约的。）
\end{enumerate}

证明 z 的中心化子与 x 的中心化子相同；它是由 x 生成的循环子群。

\begin{enumerate}
\def\labelenumi{\arabic{enumi}.}
\setcounter{enumi}{26}
\tightlist
\item
  设 \(G = \times_{A} G_{i}\) 为族 \((G_{i})_{i \in I}\) 的群沿公共子群 A 的融合和。对所有 \(i \in I\) ，选取一个子集 \(P_{i}\) ，含于 \(G_{i}\) 且满足第 3 号的条件 (A)。
\end{enumerate}

\begin{enumerate}
\def\labelenumi{(\alph{enumi})}
\tightlist
\item
  设 \(x \in G\) ，并设 \((a; i_1, \ldots, i_n; p_1, \ldots, p_n)\) 为 \(x\) 的一个既约分解；则
\end{enumerate}

\[
x = a \prod_ {\alpha = 1} ^ {n} p _ {\alpha}, \quad \text { with } \quad p _ {\alpha} \in \mathrm{P} _ {i _ {\alpha}} - \{e _ {i _ {\alpha}} \}, i _ {\alpha} \neq i _ {\alpha + 1}.
\]

证明，如果 \(i_1 \neq i_n\) ，则 \(x\) 的阶是无限的。证明，如果 \(i_1 = i_n\) ，则 \(x\) 共轭于一个其分解的长度 \(\leqslant n - 1\) 的元素。

\begin{enumerate}
\def\labelenumi{(\alph{enumi})}
\setcounter{enumi}{1}
\tightlist
\item
  由此推出，G 的每个有限阶元素都共轭于诸 \(G_{i}\) 之一的一个元素。特别地，如果诸 \(G_{i}\) 除 e 外没有有限阶元素，则 G 亦然。
\end{enumerate}

\begin{enumerate}
\def\labelenumi{\arabic{enumi}.}
\setcounter{enumi}{27}
\tightlist
\item
  我们保留前一练习的记号。
\end{enumerate}

\begin{enumerate}
\def\labelenumi{(\alph{enumi})}
\item
  设 \(\mathbf{N}\) 为 \(\mathbf{A}\) 的一个子群。假设对所有 \(i \in \mathbf{I}\) ，群 \(\mathbf{N}\) 在 \(\mathbf{G}_i\) 中正规。证明 \(\mathbf{N}\) 在 \(\mathbf{G} = \times_{\mathbf{A}} \mathbf{G}_i\) 中正规，且 \(\times_{\mathbf{A/N}} \mathbf{G}_i / \mathbf{N}\) 到 \(\mathbf{G}/\mathbf{N}\) 中的典范同态是一个同构。
\item
  对于所有 \(i \in \mathbf{I}\) ，设 \(\mathbf{H}_i\) 为 \(\mathbf{G}_i\) 的一个子群，并设 \(\mathbf{B}\) 为 \(\mathbf{A}\) 的一个子群。假设 \(\mathbf{H}_i \cap \mathbf{A} = \mathbf{B}\) 对所有 \(i\) 成立。证明 \(\star_B \mathbf{H}_i\) 到 \(\mathbf{G}\) 中的典范同态是单射；其像是 \(\mathbf{G}\) 的由诸 \(\mathbf{H}_i\) 生成的子群。
\end{enumerate}

¶ 29. 设 \(G = G_1 *_A G_2\) 为两个群 \(G_1\) 与 \(G_2\) 沿子群 A 的融合和。

\begin{enumerate}
\def\labelenumi{(\alph{enumi})}
\item
  证明若 \(G_{1}\) 与 \(G_{2}\) 都是有限生成的，则 G 是有限生成的。
\item
  假设 \(\mathbf{G}_1\) 与 \(\mathbf{G}_2\) 是有限表现的。证明下列两个性质的等价性：
\end{enumerate}

\begin{enumerate}
\def\labelenumi{(\arabic{enumi})}
\item
  G 是有限表现的。
\item
  A 是有限生成的。
\end{enumerate}

（首先证明 (1) 蕴含 (2)。另一方面，若 A 不是有限生成的，则存在 A 的子群的一个右有向族 \((\mathbf{A}_{i})_{i\in I}\) ，满足 \(\bigcup_{i\in I}\mathbf{A}_{i}=\mathbf{A}\) 且对所有 i 有 \(A_{i}\neq A\) 。设 \(H_{i}\) （相应地 H）为 \(G_{1}*G_{2}\) 到 \(G_{1}*_{A_{i}}G_{2}\) （相应地到 G）上的典范同态的核。证明 H 是诸 \(H_{i}\) 的并，且对所有 i 有 \(H_{i}\neq H\) ；然后应用习题 16(d)。）

\begin{enumerate}
\def\labelenumi{(\alph{enumi})}
\setcounter{enumi}{2}
\tightlist
\item
  由此给出一个有限生成但非有限表现的群的例子。（取 \(G_{1}\) 与 \(G_{2}\) 为秩 2 的自由群，A 为有限秩的自由群，参见习题 22。）
\end{enumerate}

\begin{enumerate}
\def\labelenumi{\arabic{enumi}.}
\setcounter{enumi}{29}
\tightlist
\item
  设 A 和 B 是两个群，G = A * B 为它们的自由积。
\end{enumerate}

\begin{enumerate}
\def\labelenumi{(\alph{enumi})}
\tightlist
\item
  设 \(\mathbf{Z}\) 为 \(\mathbf{A} \backslash \mathbf{G} / \mathbf{A}\) 中异于 \(\mathbf{A}\) 的一个元素。证明 \(\mathbf{Z}\) 包含且仅包含一个如下形式的元素 \(z\) ：
\end{enumerate}

\[
b _ {1} a _ {1} b _ {2} a _ {2} \dots b _ {n - 1} a _ {n - 1} b _ {n},
\]

其中 \(n \geqslant 1\) , \(a_i \in \mathbf{A} - \{e\}\) , \(b_i \in \mathbf{B} - \{e\}\) 。证明 \(\mathbf{Z}\) 的每个元素都可以唯一地写成形式 \(aza'\) ，其中 \(a, a' \in \mathbf{A}\) 。

\begin{enumerate}
\def\labelenumi{(\alph{enumi})}
\setcounter{enumi}{1}
\tightlist
\item
  设 \(x \in \mathbf{G} - \mathbf{A}\) 。设 \(f_x\) 为 \(\mathbf{A} * \mathbf{A}\) 到 \(\mathbf{G}\) 中的同态，它在 \(\mathbf{A} * \mathbf{A}\) 的第一个（相应地，第二个）因子上的限制是恒同映射（相应地，映射 \(a \mapsto xax^{-1}\) ）。证明 \(f_x\) 是单射。（将 \(x\) 写成上述形式 \(aza'\) ，从而化归为 \(x = z\) 的情形；利用 \(\mathbf{G}\) 中既约分解的唯一性。）
\end{enumerate}

特别地， \(\mathbf{A} \cap x\mathbf{A}x^{-1} = \{e\}\) ，且 \(\mathbf{A}\) 在 \(\mathbf{G}\) 中的正规化子是 \(\mathbf{A}\) 。

\begin{enumerate}
\def\labelenumi{(\alph{enumi})}
\setcounter{enumi}{2}
\tightlist
\item
  设 \(\mathbf{T}\) 为 \(\mathbf{A} \backslash \mathbf{G} / \mathbf{B}\) 的一个元素。证明 \(\mathbf{T}\) 包含且仅包含一个如下形式的元素 \(t\) ：
\end{enumerate}

\[
b _ {1} a _ {1} b _ {2} a _ {2} \dots b _ {n - 1} a _ {n - 1},
\]

其中 \(n \geqslant 1\) , \(a_i \in \mathbf{A} - \{e\}\) , \(b_i \in \mathbf{B} - \{e\}\) 。证明 \(\mathbf{T}\) 的每个元素都可以唯一地写成 atb 的形式，其中 \(a \in \mathbf{A}\) , \(b \in \mathbf{B}\) 。

\begin{enumerate}
\def\labelenumi{(\alph{enumi})}
\setcounter{enumi}{3}
\tightlist
\item
  设 \(x \in G\) ，并设 \(h_x\) 为 \(G\) 的唯一自同态，使得 \(h_x(a) = a\) 当 \(a \in A\) ，且 \(h_x(b) = x b x^{-1}\) 当 \(b \in B\) 。证明 \(h_x\) 是单射。（用与 (b) 中相同的方法，利用 (c) 中的分解 \(x = a t b\) 。）特别地， \(A \cap x B x^{-1} = \{e\}\) 。
\end{enumerate}

给出一个 \(h_{x}\) 不是满射的例子。

\begin{enumerate}
\def\labelenumi{\arabic{enumi}.}
\setcounter{enumi}{30}
\tightlist
\item
  设 G 为某个群的族 \((G_{i})_{i\in I}\) 的自由积。设 S 为 G 的与诸 \(G_{i}\) 之一共轭的子群组成的集合，并设 \(\Theta\) 为 S 中所有元素的并。
\end{enumerate}

\begin{enumerate}
\def\labelenumi{(\alph{enumi})}
\item
  设 \(\mathbf{H}, \mathbf{H}' \in \mathbf{S}\) ，其中 \(\mathbf{H} \neq \mathbf{H}'\) ，并设 \(x \in \mathbf{H} - \{e\}\) , \(x' \in \mathbf{H}' - \{e\}\) 。证明 \(xx' \in \mathbf{G} - \Theta\) 。（化归到 \(\mathbf{H}\) 是诸 \(\mathbf{G}_i\) 之一的情形；将 \(\mathbf{H}'\) 写成形式 \(y\mathrm{G}_j y^{-1}\) ，并像前一习题那样进行。）
\item
  设 \(\mathbf{C}\) 为 \(\mathbf{G}\) 的一个含于 \(\Theta\) 的子群。证明存在 \(\mathbf{H} \in \mathbf{S}\) ，使得 \(\mathbf{C} \subset \mathbf{H}\) 。
\item
  设 C 为 G 的一个其所有元素均为有限阶的子群。证明 C 含于 \(\Theta\) （参见习题 27）；由此推出 C 含于诸 \(G_{i}\) 之一的一个共轭之中。
\end{enumerate}

¶ 32. 设 A 和 B 是两个群，G = A * B 是它们的自由积。设 X 为形如 \((a, b)\) 、其中 \(a \in A - \{e\}\) 且 \(b \in B - \{e\}\) 的换位子组成的集合，并设 \(S = X \cup X^{-1}\) 。

\begin{enumerate}
\def\labelenumi{(\alph{enumi})}
\item
  证明 \(\mathbf{X} \cap \mathbf{X}^{-1} = \varnothing\) ，并且，若 \(x\) 是 \(\mathbf{X}\) 的一个元素，则存在唯一的有序对 \(a, b\) ，属于 \((\mathbf{A} - \{e\}) \times (\mathbf{B} - \{e\})\) ，使得 \((a, b) = x\) 。
\item
  设 \(n \geqslant 1\) ，并设 \(s_1, \ldots, s_n\) 为 \(S\) 中元素的一个序列，使得 \(s_i s_{i+1} \neq e\) 对 \(1 \leqslant i < n\) 成立；设 \(a_n \in A - \{e\}, b_n \in B - \{e\}, \varepsilon(n) = \pm 1\) 满足 \(s_n = (a_n, b_n)^{\varepsilon(n)}\) 。设 \(g = s_1 \ldots s_n\) 。通过对 \(n\) 作归纳证明： \(g\) 的既约分解的长度 \(\geqslant n + 3\) ，并且或者以 \(a_nb_n\) 结束（若 \(\varepsilon(n) = 1\)），或者以 \(b_na_n\) 结束（若 \(\varepsilon(n) = -1\)）。
\item
  设 R 为典范投影 \(A*B \rightarrow A \times B\) 的核，该投影（第 3 号）对应于典范单射 \(A \rightarrow A \times B\) 与 \(B \rightarrow A \times B\) 。证明 R 是以 X 为基本族的自由群。（利用 (b) 证明 X 是自由的；若 \(R_{X}\) 为 G 中由 X 生成的子群，证明 \(R_{X}\) 是正规的；由此推出它与 R 重合。）
\end{enumerate}

\begin{enumerate}
\def\labelenumi{\arabic{enumi}.}
\setcounter{enumi}{32}
\tightlist
\item
  设 \(E = G *_{A} G'\) 为两个群沿公共子群 A 的融合和。设 P（相应地 \(P'\) ）为 G（相应地 \(G'\) ）的满足第 3 号的条件 (A) 的子集。
\end{enumerate}

设 \(n\) 为一个整数 \(\geqslant 1\) 。设 \(L_n\) 表示 \(x \in E\) 中既约分解的长度 \(\leqslant 2n - 1\) 者组成的集合。设 \(\mathbf{M}_n\) 表示 \(x \in \mathbf{E}\) 中形如 \(ap_1p_1'p_2p_2'\ldots p_np_n'\) 的元素组成的集合，其中 \(a \in \mathbf{A}, p_i \in \mathbf{P} - \{e\}, p_i' \in \mathbf{P}' - \{e\}\) ；类似地，令 \(\mathbf{M}_n'\) 表示形如 \(ap_1'p_1p_2'p_2\ldots p_np_n'\) 的元素组成的集合，其中 \(a, p_i, p_i'\) 满足相同的条件。

\begin{enumerate}
\def\labelenumi{(\alph{enumi})}
\item
  证明 \(L_{n}\) , \(M_{n}\) 与 \(M_{n}^{\prime}\) 两两不相交，且它们的并是 E 中既约分解的长度 \(\leqslant2n\) 的元素组成的集合。
\item
  构造一个双射 \(\varepsilon\) ——\(\mathbf{M}_n\) 到 \(\mathbf{M}_n'\) 上的双射——使得 \(\varepsilon(ax) = a\varepsilon(x)\) （若 \(a \in \mathbf{A}\) , \(x \in \mathbf{M}_n\)）。
\item
  设 \(\mathbf{X}_n = \mathbf{L}_n \cup \mathbf{M}_n\) 与 \(\mathbf{X}_n' = \mathbf{L}_n \cup \mathbf{M}_n'\) ；证明 \(\mathbf{G} \cdot \mathbf{X}_n = \mathbf{X}_n\) 且 \(\mathbf{G}' \cdot \mathbf{X}_n' = \mathbf{X}_n'\) 。
\item
  \(\varepsilon\) 扩充为 \(\mathbf{X}_n\) 到 \(\mathbf{X}_n'\) 上的一个双射，扩充方式为令 \(\varepsilon(x) = x\) （若 \(x \in \mathbf{L}_n\)）。证明 E 可以如下作用在 \(\mathbf{X}_n\) 上： \(g(x) = gx\) （若 \(g \in G\) , \(x \in \mathbf{X}_n\)），且 \(g'(x) = \varepsilon^{-1}(g' \varepsilon (x))\) （若 \(g' \in G'\) , \(x \in \mathbf{X}_n\)）。证明，若 \(g \in \mathbf{X}_n\) ，则 \(g(e) = g\) 。
\item
  设 \(\tau_{n}\colon\mathbf{E}\to\mathfrak{S}_{\mathbf{X}_{n}}\) 为由上述 \(\mathbf{E}\) 在 \(\mathbf{X}_{n}\) 上的作用所定义的同态，并令 \(\mathbf{R}_{n}\) 为 \(\tau_{n}\) 的核。证明 \(\mathbf{R}_{n}\cap\mathbf{X}_{n}=\{e\}\) ；由此推出诸 \(\mathbf{R}_{n}\) 的交化为 \(e\) 。
\item
  当 G 与 G' 有限时，证明 E 是剩余有限的（参见 § 5 习题 5）。（注意此时诸 R \(_{n}\) 都是有限指数的子群。）
\end{enumerate}

¶ 34. (a) 设 \((\mathrm{H}_{i})\) 为群 G 的一个右有向正规子群族。假设 \(\bigcap_{i} \mathrm{H}_{i} = \{e\}\) 。证明：对每个群 \(\mathbf{G}'\) ，典范同态 \(\mathbf{G}' * \mathbf{G} \to \mathbf{G}' * (\mathbf{G}/\mathbf{H}_{i})\) 的核的交等于 \(\{e\}\) 。

\begin{enumerate}
\def\labelenumi{(\alph{enumi})}
\setcounter{enumi}{1}
\item
  设 G 为某个群的族 \((G_{\alpha})\) 的自由积。证明，如果所有 \(G_{\alpha}\) 都是剩余有限的（参见 § 5 习题 5），那么 G 也是。（首先化归为族有限的情形，然后化归为族有两个元素的情形；由 (a)，可设诸 \(G_{\alpha}\) 有限，再应用前一习题。）
\item
  由此推出，每个自由群都是剩余有限的。
\end{enumerate}

\begin{enumerate}
\def\labelenumi{\arabic{enumi}.}
\setcounter{enumi}{34}
\item
  *设 G（相应地，G'）为形如 \(a/b\) 、其中 \(a \in \mathbf{Z}\) , \(b \in \mathbf{Z}\) 且 \(b\) 不被 2（相应地，不被 3）整除的分数组成的加法群。设 E = G * \(_{\mathbf{Z}}\) G' 为 G 与 G' 沿其公共子群 Z 的融合和。证明 G 与 G' 是剩余有限的，但 E 不是。（证明 E 的有限指数子群只有 E 本身。）*
\item
  设 \((\mathbf{G}_1, \ldots, \mathbf{G}_n)\) 与 \((\mathbf{A}_1, \ldots, \mathbf{A}_{n-1})\) 为两个群列；对每个 \(i\) ，设 \(\phi_i: \mathbf{A}_i \to \mathbf{G}_i\) 与 \(\psi_i: \mathbf{A}_i \to \mathbf{G}_{i+1}\) 为两个单同态；通过这些同态，\(\mathbf{A}_i\) 既可与 \(\mathbf{G}_i\) 的一个子群等同，也可与 \(\mathbf{G}_{i+1}\) 的一个子群等同。
\end{enumerate}

\begin{enumerate}
\def\labelenumi{(\alph{enumi})}
\item
  设 \(H_{1}=G_{1}, H_{2}=H_{1} *_{A_{1}} G_{2}, \ldots, H_{n}=H_{n-1} *_{A_{n-1}} G_{n}\) 。群 \(H_{n}\) 称为群 \((G_{i})\) 沿子群 \((A_{i})\) 的融合和；记为 \(G_{1} *_{A_{1}} G_{2} *_{A_{2}} \cdots *_{A_{n-1}} G_{n}\) 。若 T 为任意群，在 \(H_{n}\) 到 T 中的同态与同态族 \((t_{1},\ldots ,t_{n})\) （其中 \(t_i\colon \mathbf{G}_i\to \mathbf{T}\)）之间定义一个双射，使得 \(t_i\circ \phi_i = t_{i + 1}\circ \psi_i\) 对 \(1\leqslant i\leqslant n - 1\) 成立。
\item
  将每个 \(G_{i}\) 与 \(H_{n}\) 中相应的子群等同。证明 \(G_{i} \cap G_{i+1} = A_{i}\) ，并且 \(G_{i} \cap G_{i+2} = A_{i} \cap A_{i+1}\) （后一交集取在 \(G_{i+1}\) 中）。
\end{enumerate}

¶ 37. 设 G 为一个群，S 为 G 的有限对称子集且生成 G。令 \(\mathfrak{T}\) 表示 G 的有限子集全体，按包含关系排序；它是一个有向集。

\begin{enumerate}
\def\labelenumi{(\alph{enumi})}
\item
  设 \(\mathbf{T} \in \mathfrak{T}\) ，并设 \(\mathbf{B}_{\mathbf{T}}\) 为 \(\mathbf{G} - \mathbf{T}\) 的连通 S-分量组成的集合（参见习题 18）；设 \(\mathbf{B}_{\mathbf{T}}^{\infty} = \mathbf{B}_{\mathbf{T}} - (\mathfrak{T} \cap \mathbf{B}_{\mathbf{T}})\) 。证明 \(\mathbf{B}_{\mathbf{T}}\) 是有限的（注意，若 \(\mathbf{T} \neq \varnothing\) 且 \(\mathbf{Z} \in \mathbf{B}_{\mathbf{T}}\) ，则存在 \(t \in \mathbf{T}\) 与 \(z \in \mathbf{Z}\) ，它们是 S-相邻的）。
\item
  设 \(\mathbf{T}, \mathbf{T}' \in \mathfrak{T}\) 满足 \(\mathbf{T} \subset \mathbf{T}'\) 。若 \(Z' \in \mathbf{B}_{\mathbf{T}'}\) ，证明存在唯一的元素 \(Z \in \mathbf{B}_{\mathbf{T}}\) 使得 \(Z' \subset Z\) ；映射 \(f_{\mathbf{T}\mathbf{T}'}: \mathbf{B}_{\mathbf{T}'} \to \mathbf{B}_{\mathbf{T}}\) （满足 \(f_{\mathbf{T}\mathbf{T}'}(\mathbf{Z}') = \mathbf{Z}\)）把 \(\mathbf{B}_{\mathbf{T}}^{\infty}\) 映到 \(\mathbf{B}_{\mathbf{T}}^{\infty}\) 上。诸 \(\mathbf{B}_{\mathbf{T}}\) （其中 \(\mathbf{T}\) 遍历 \(\mathfrak{T}\)）的逆向极限记作 \(\mathbf{B}\) 。证明 \(B = \varprojlim B_{\mathbf{T}}^{\infty}\) ，且 \(B\) 在 \(B_{\mathbf{T}}\) 中的像是 \(B_{\mathbf{T}}^{\infty}\) 。集合 \(B\) 称为 \(G\) 的端集；它非空当且仅当 \(G\) 是无限的。
\item
  设 \(\mathbf{T}\) 为 \(\mathbf{G}\) 的一个非空 S-连通有限子集，并设 \(\mathbf{Z} \in \mathbf{B}_{\mathbf{T}}^{\infty}\) 。证明存在 \(g \in \mathbf{G}\) ，使得 \(g\mathbf{T} \cap \mathbf{T} = \varnothing\) 且 \(g\mathbf{T} \cap \mathbf{Z} \neq \varnothing\) ，而此时 \(g\mathbf{T} \subset \mathbf{Z}\) 。证明存在 \(\mathbf{Z}_1 \in \mathbf{B}_{\mathbf{T}}\) ，使得 \(g\mathbf{Z}_1\) 包含 \(\mathbf{T}\) 以及所有的 \(\mathbf{Z}' \in \mathbf{B}_{\mathbf{T}}\) （满足 \(\mathbf{Z}' \neq \mathbf{Z}\) ）（注意 \(\mathbf{T} \cup \bigcup_{\mathbf{Z}' \neq \mathbf{Z}} \mathbf{Z}'\) 是 S-连通的且不与 \(g\mathbf{T}\) 相交）；由此推出，若 \(\mathbf{Z}' \in \mathbf{B}_{\mathbf{T}}\) 不同于 \(\mathbf{Z}_1\) ，则 \(g\mathbf{Z}' \subset \mathbf{Z}\) 。证明，记 \(\mathbf{T}' = \mathbf{T} \cup g\mathbf{T}\) ，则 \(\mathbf{Z}\) 在 \(\mathbf{B}_{\mathbf{T}}^{\infty}\) 中的原像至少由 \(n - 1\) 个元素组成，其中 \(n = \operatorname{Card}(\mathbf{B}_{\mathbf{T}}^{\infty})\) 。
\item
  设 \(S'\) 为 \(G\) 的一个有限对称子集，它生成 \(G\) ，并设 \(B'\) 为 \(G\) 的、相对于 \(S'\) 的端集。定义 \(B\) 到 \(B'\) 上的一个双射（化归为 \(S \subset S'\) 的情形）。于是 \(B\) 的基数与 \(S\) 的选取无关；它称为 \(G\) 的端数。
\item
  利用 (c) 证明，G 的端数等于 0、1、2 或 \(2^{\aleph_0}\) 。
\item
  证明 \(\mathbf{Z}\) 的端数是 2，而 \(\mathbf{Z}^n\) ( \(n \geqslant 2\) ) 的端数是 1。
\end{enumerate}

¶ 38. 设 X 为有限集，F(X) 为在 X 上构造的自由群，且 S = X ∪ X \(^{-1}\) 。

\begin{enumerate}
\def\labelenumi{(\alph{enumi})}
\item
  设 \(\mathbf{S}_n\) 为乘积 \(s_1 \ldots s_m\) , \(s_i \in \mathbf{S}\) , \(m \leqslant n\) 组成的集合。证明 \(\mathrm{F}(\mathbf{X}) - \mathbf{S}_n\) 的每个连通 S-分量（参见习题 18）包含且仅包含一个形如 \(s_1 \ldots s_{n+1}\) 的元素，其中 \(s_i \in \mathbf{S}\) 且 \(s_i s_{i+1} \neq e\) 对 \(1 \leqslant i \leqslant n\) 成立。
\item
  由 (a) 推出集合 B（即 F(X) 的端集，参见习题 37）到无穷序列 \((s_{1},\ldots,s_{n},\ldots)\) （其元素属于 S 且满足 \(s_{i}s_{i+1}\neq e\) 对所有 i 成立）组成的集合上的一个双射。特别地，F(X) 的端数等于 \(2^{\aleph_0}\) （若 \(\text{Card}(X)\geqslant2\)）。
\end{enumerate}

\begin{enumerate}
\def\labelenumi{\arabic{enumi}.}
\setcounter{enumi}{38}
\tightlist
\item
  *设 G 是一个有限生成群，F 是 G 的包含 e 且生成 G 的对称有限子集的集合。
\end{enumerate}

\begin{enumerate}
\def\labelenumi{(\alph{enumi})}
\tightlist
\item
  若 \(\mathbf{X} \in \mathfrak{F}\) 且 \(n \in \mathbf{N}\) ，设 \(d_n(\mathbf{X})\) 表示 \(\mathbf{G}\) 中形如 \(x_1 \ldots x_n\) （其中 \(x_i \in \mathbf{X}\)）的元素的个数。若 \(\mathbf{X}, \mathbf{Y} \in \mathfrak{F}\) ，证明存在一个整数 \(a \geqslant 1\) ，使得
\end{enumerate}

\[
d _ {n} (\mathbf {X}) \leqslant d _ {a n} (\mathbf {Y}) \quad \text { for   all } n \in \mathbf {N}.
\]

由此推出 \(\limsup_{n\to \infty}\frac{\log(d_n(\mathbf{X}))}{\log(n)}\) 不依赖于 \(\mathbf{X}\) 在 \(\mathfrak{F}\) 中的选取；这一极限（有限或无限）记作 \(e(\mathbf{G})\) 。

\begin{enumerate}
\def\labelenumi{(\alph{enumi})}
\setcounter{enumi}{1}
\item
  若 \(\mathbf{X} \in \mathfrak{F}\) 且 \(\mathbf{G}\) 是无限的，证明 \(d_n(\mathbf{X}) < d_{n+1}(\mathbf{X})\) 对所有 \(n\) 成立。由此推出 \(e(\mathbf{G})\) \(\geqslant 1\) 。
\item
  设 H 为一个有限生成群。证明 \(e(\mathbf{H}) \leqslant e(\mathbf{G})\) ，如果 H 同构于 G 的一个子群（相应地，一个商群）。若 E 是 G 被 H 的扩张，证明 \(e(\mathbf{E}) \geqslant e(\mathbf{H}) + e(\mathbf{G})\) ，且当 \(E = G \times H\) 时等号成立。
\item
  证明 \(e(\mathbf{Z}^n) = n\) 。
\item
  设 \(\mathbf{X} \in \mathfrak{F}\) 。考虑下列性质：
\item
  存在一个实数 \(a > 1\) ，使得 \(d_n(\mathbf{X}) \geqslant a^n\) 对所有充分大的 \(n\) 成立。
\end{enumerate}

证明：若性质 (i) 对 \(\mathfrak{F}\) 的某个元素成立，则它对 \(\mathfrak{F}\) 的每个元素都成立。此时称 \(\mathbf{G}\) 为指数型的；这蕴含 \(e(\mathbf{G}) = +\infty\) 。

\begin{enumerate}
\def\labelenumi{(\alph{enumi})}
\setcounter{enumi}{5}
\item
  设 \(G_{1}\) 与 \(G_{2}\) 为包含同一子群 A 的两个群，令 \(H = G_{1} *_{A} G_{2}\) 为相应的融合和。假设 \(G_{1}\) 与 \(G_{2}\) 是有限生成的（从而 H 也是），且对 i = 1, 2 有 \(G_{i} \neq A\) 。证明：若指标 ( \(G_{1}:A\) ), ( \(G_{2}:A\) ) 中至少有一个 \(\geqslant 3\) ，则 H 是指数型的。
\item
  证明秩 \(\geqslant 2\) 的自由群是指数型的。*
\end{enumerate}

\begin{enumerate}
\def\labelenumi{\arabic{enumi}.}
\setcounter{enumi}{39}
\tightlist
\item
  *设 \(n\) 为一个整数 \(\geqslant 2\) ，\(T_n\) 为对角项都等于 1 的 \(n\) 阶上三角矩阵（在 \(Z\) 上）组成的群。设 \(e(T_n)\) 为习题 39 中定义的群 \(T_n\) 的不变量。证明
\end{enumerate}

\[
e (\mathrm{T} _ {n}) \leqslant \sum_ {i = 1} ^ {n} i (n - i). *
\]

\BourbakiExercisesSection

\begin{enumerate}
\def\labelenumi{\arabic{enumi}.}
\item
  确定一个具有 \(n\) 个元素的集合上的所有环结构，其中 \(2 \leqslant n \leqslant 7\) ，并确定这些环的理想。
\item
  设 A 是一个环。
\end{enumerate}

\begin{enumerate}
\def\labelenumi{(\alph{enumi})}
\item
  将 \(a \in A\) 对应于左位似 \(x \mapsto ax\) 的映射是 A 到 A 的加法群的与诸右位似交换的自同态环上的一个同构。
\item
  设 M 为 A 的四个元素组成的序列 \((a, b, c, d)\) 的集合，这些序列可以表示成表格或矩阵 \(\begin{pmatrix} a & b \\ c & d \end{pmatrix}\) 的形式。对 M 的每个元素 \(\begin{pmatrix} a & b \\ c & d \end{pmatrix}\) ，我们使映射 \((x, y) \mapsto (ax + by, cx + dy)\) 与之对应，它是 \(\mathbf{A} \times \mathbf{A}\)
\end{enumerate}

到其自身的映射。证明由此得到 M 到与运算 \((x, y) \mapsto (xr, yr)\) （对所有 \(r \in A\)）交换的交换群 A × A 的自同态环 B 上的一个双射。此后将 M 与环 B 等同。计算两个矩阵的乘积。

\begin{enumerate}
\def\labelenumi{(\alph{enumi})}
\setcounter{enumi}{2}
\tightlist
\item
  证明子集 \(\mathbf{T}\) ——由 \(\mathbf{M}\) 中形如 \(\begin{pmatrix} a & b \\ 0 & c \end{pmatrix}\) 的矩阵组成——是 \(\mathbf{M}\) 的一个使子群 \(\mathbf{A} \times 0\) 在 \(\mathbf{A} \times \mathbf{A}\) 中保持稳定的子环。
\end{enumerate}

映射 \(\begin{pmatrix} a & b \\ 0 & c \end{pmatrix} \mapsto a\) 与 \(\begin{pmatrix} a & b \\ 0 & c \end{pmatrix} \mapsto c\) 是 T 到 A 中的同态；设 \(a\) 与 \(b\) 为它们的核。证明 \(a + b = T\) ，\(ab = \{0\}\) ，且 \(ba = b \cap a\) \(\neq \{0\}\) （若 \(A \neq \{0\}\) ；参见第 9 号，命题 6）。

\begin{enumerate}
\def\labelenumi{(\alph{enumi})}
\setcounter{enumi}{3}
\tightlist
\item
  若 \(\mathbf{A} = \mathbf{Z} / 2\mathbf{Z}\) ，则 \(\mathbf{T}\) 不是交换的，但 \(\mathbf{T}\) 的可逆元素的乘法群是交换的。
\end{enumerate}

\begin{enumerate}
\def\labelenumi{\arabic{enumi}.}
\setcounter{enumi}{2}
\item
  设 A 为一个环，且 \(a \in A\) 。如果存在且仅存在一个 \(a' \in A\) 使得 \(aa' = 1\) ，则 a 是可逆的，并且 \(a' = a^{-1}\) （首先证明 a 是左可消去的，然后考虑乘积 \(aa'a\)）。
\item
  设 A 为一个伪环，其中 A 的每个加法子群都是 A 的左理想。
\end{enumerate}

\begin{enumerate}
\def\labelenumi{(\alph{enumi})}
\item
  若 \(\mathbf{A}\) 是一个环，则 \(\mathbf{A}\) 典范同构于 \(\mathbf{Z} / n\mathbf{Z}\) ，其中 \(n\) 为某个适当的整数。
\item
  如果 A 没有单位元且每个非零元素都是可消去的，则 A 同构于 Z 的一个伪子环。
\end{enumerate}

\begin{enumerate}
\def\labelenumi{\arabic{enumi}.}
\setcounter{enumi}{4}
\tightlist
\item
  设 \(\mathbf{M}\) 为一个幺半群，\(\mathbf{Z}^{(\mathbf{M})}\) 为以 \(\mathbf{M}\) 为基的自由交换群，\(u\colon \mathbf{M}\to \mathbf{Z}^{(\mathbf{M})}\) 为典范映射。
\end{enumerate}

\begin{enumerate}
\def\labelenumi{(\alph{enumi})}
\item
  证明在 \(\mathbf{Z}^{(\mathbf{M})}\) 上存在唯一的乘法，使得 \(\mathbf{Z}^{(\mathbf{M})}\) 是一个环，且 \(u\) 是幺半群 \(\mathbf{M}\) 到环 \(\mathbf{Z}^{(\mathbf{M})}\) 的乘法幺半群中的幺半群态射。
\item
  设 A 为一个环，\(v: \mathbf{M} \to \mathbf{A}\) 为幺半群 M 到 A 的乘法幺半群中的幺半群态射。存在唯一的环态射 \(f: \mathbf{Z}^{(\mathbf{M})} \to \mathbf{A}\) ，使得
\end{enumerate}

\[
f \circ u = v.
\]

\begin{enumerate}
\def\labelenumi{\arabic{enumi}.}
\setcounter{enumi}{5}
\tightlist
\item
  设 A 为一个交换环，M 为 A 的加法群的一个子幺半群。设 n 为一个整数 \textgreater0 ，使得 \(m^{n} = 0\) 对 \(m \in M\) 成立。
\end{enumerate}

\begin{enumerate}
\def\labelenumi{(\alph{enumi})}
\item
  若 \(n!\) 在 \(\mathbf{A}\) 中不是零因子，则由 \(n\) 个 \(\mathbf{M}\) 中元素组成的族之乘积为零。
\item
  证明：若 \(n!\) 在 A 中是零因子，则 (a) 中的结论不一定成立。
\end{enumerate}

\begin{enumerate}
\def\labelenumi{\arabic{enumi}.}
\setcounter{enumi}{6}
\tightlist
\item
  对于所有 \(n \in \mathbf{N}^*\) ,
\end{enumerate}

\[
n! = \sum_ {\nu = 0} ^ {n - 1} (- 1) ^ {\nu} \binom {n} {\nu} (n - \nu) ^ {n}
\]

（利用第 2 号，命题 2。）

\begin{enumerate}
\def\labelenumi{\arabic{enumi}.}
\setcounter{enumi}{7}
\item
  在环中，由一个左理想生成的右理想是一个双边理想。
\item
  在环中，一个右理想的右零化子是一个双边理想。
\item
  在环 A 中，由元素 xy - yx（其中 x 与 y 遍历 A）生成的双边理想，是使 A/a 为交换环的双边理想 a 中最小的一个。
\end{enumerate}

¶ 11. 在环 A 中，一个双边理想 a 称为不可约的，如果不存在有序对 (b, c)，其中 b、c 是与 a 不同的双边理想，且满足 a = b ∩ c。

\begin{enumerate}
\def\labelenumi{(\alph{enumi})}
\item
  证明 A 的所有不可约双边理想的交化为 0（注意，不含任何元素 \(a \neq 0\) 的双边理想组成的集合是归纳的，并应用佐恩引理）。
\item
  由此推出，每个双边理想 \(\mathfrak{a}\) （\(\mathbf{A}\) 的）是包含它的所有不可约理想的交。
\end{enumerate}

\begin{enumerate}
\def\labelenumi{\arabic{enumi}.}
\setcounter{enumi}{11}
\tightlist
\item
  设 A 为一个环，且 \((\mathfrak{a}_{i})_{i\in I}\) 为 A 的左理想的有限族，使得 A 的加法群是诸加法群 \(a_{i}\) 的直和。若 \(1=\sum_{i\in I}e_{i}(e_{i}\in\mathfrak{a}_{i})\) ，则 \(e_{i}^{2}=e_{i}, e_{i}e_{j}=0\) （若 \(i\neq j\)）且 \(a_{i}=Ae_{i}\) （对所有 \(x\in A\) 写出 x=x.1）。反之，如果 \((e_{i})_{i\in I}\) 是 A 的幂等元的有限族，满足 \(e_{i}e_{j}=0\) （对 \(i\neq j\)）且 \(1=\sum_{i}e_{i}\) ，则 A 是诸左理想 \(Ae_{i}\) 的直和。为使诸理想 \(Ae_{i}\) 是双边的，必要且充分条件是诸 \(e_{i}\) 属于 A 的中心。
\end{enumerate}

¶ 13. 设 A 是一个环，e 是 A 的一个幂等元。

\begin{enumerate}
\def\labelenumi{(\alph{enumi})}
\item
  证明 \(\mathbf{A}\) 的加法群是左理想 \(\mathfrak{a} = \mathbf{A}e\) 与左零化子 \(\mathfrak{b}\) （对 \(e\) 而言）的直和（注意，对所有 \(x\in \mathbf{A}\) ，有 \(x - xe\in \mathfrak{b}\)）。
\item
  每个右理想 \(\mathfrak{d}\) （属于 \(\mathbf{A}\)）是 \(\mathfrak{d} \cap \mathfrak{a}\) 与 \(\mathfrak{d} \cap \mathfrak{b}\) 的直和。
\item
  若 \(\mathbf{Ae} = e\mathbf{A}\) ，则 \(\mathfrak{a}\) 与 \(\mathfrak{b}\) 是 \(\mathbf{A}\) 的双边理想，并且定义 \(\mathbf{A}\) 的一个直分解。
\end{enumerate}

\begin{enumerate}
\def\labelenumi{\arabic{enumi}.}
\setcounter{enumi}{13}
\item
  设 A 为一个交换环，其中零因子只有有限 n 个。证明 A 至多有 \((n+1)^{2}\) 个元素。（设 \(a_{1},\ldots,a_{n}\) 为零因子。证明零化子 \(J_{1}\) （\(a_{1}\) 的）至多有 \(n+1\) 个元素，并且商环 A/ \(J_{1}\) 至多有 \(n+1\) 个元素。）
\item
  环胚是一个带有两个合成律的集合 E：（a）一个结合乘法 \(xy\) ；（b）一个以加法记、并非处处有定义的合成律（§ 1，第 1 号），且满足下列条件：
\end{enumerate}

\begin{enumerate}
\def\labelenumi{(\arabic{enumi})}
\item
  它是交换的（换句话说，若 \(x + y\) 有定义，则 \(y + x\) 亦有定义且 \(x + y = y + x\) ；此时称 \(x\) 与 \(y\) 是可相加的）；
\item
  若 \(x\) 与 \(y\) 可相加，则为使 \(x + y\) 与 \(z\) 可相加，只需 \(x\) 与 \(y\) 、以及 \(y\) 与 \(z\) 分别可相加；此时 \(x\) 与 \(y + z\) 可相加，且 \((x + y) + z = x + (y + z)\) ；
\item
  存在一个单位元 0；
\item
  若 \(x\) 与 \(z\) 、以及 \(y\) 与 \(z\) 分别可相加，且 \(x + z = y + z\) ，则 \(x = y\) ；
\item
  若 \(x, y\) 可相加，则 \(xz\) 与 \(yz\) （相应地 \(zx\) 与 \(zy\)）亦可相加，且
\end{enumerate}

\[
(x + y) z = x z + y z
\]

（相应地 \(z(x + y) = zx + zy)\) 对所有 \(z \in \mathbf{E}\) 成立。

每个环都是一个环状结构。

\begin{enumerate}
\def\labelenumi{(\alph{enumi})}
\item
  考察 § 8 的定义和结果以及上述习题如何推广到环胚（环胚 E 的左理想是 E 的一个子集 \(\mathfrak{a}\) ，它在加法下稳定，且满足 \(\mathbf{E}.\mathfrak{a}\subset \mathfrak{a}\)）。
\item
  设 G 为一个带算子的群，f 与 g 为 G 的两个自同态。为使映射 \(x \mapsto f(x)g(x)\) 是 G 的自同态，必要且充分条件是子群 \(f(\mathrm{G})\) 的每个元素与 \(g(\mathrm{G})\) 的每个元素可交换；若把这一自同态记作 \(f + g\) ，且 fg 为复合自同态 \(x \mapsto f(g(x))\) ，证明 G 的具有这些合成律的自同态集合 E 是一个有单位元的环胚；为使 E 成为一个环，必要且充分条件是 G 是交换的。
\end{enumerate}

一个元素 \(f \in E\) 能与 E 的所有元素相加，必要且充分条件是 \(f(G)\) 含于 G 的中心；这些自同态组成的集合 N 是一个伪环，称为环胚 E 的核。

一个自同态 \(f\) （\(G\) 的）称为正规的，如果它与 \(G\) 的所有内自同构可交换；对每个正规稳定子群 \(H\) ，\(G, f(G)\) 都是 \(G\) 的一个正规稳定子群。证明 \(D\) ——\(G\) 的正规自同态组成的集合——构成 \(E\) 的一个子环胚，且核 \(N\) 是一个双边理想，位于 \(D.†\)

\begin{enumerate}
\def\labelenumi{\arabic{enumi}.}
\setcounter{enumi}{15}
\tightlist
\item
  给出理想 \(a, b, c\) （在环 \(\mathbf{Z}\) 中）的例子，使得 \(ab \neq a \cap b\) 且 \((a + b)(a + c) \neq a + bc\) 。
\end{enumerate}

\BourbakiExercisesSection

\begin{enumerate}
\def\labelenumi{\arabic{enumi}.}
\item
  在 § 8 习题 1 所确定的环结构中，哪些是域结构？
\item
  每个非零元素都可消去的有限环是域（参见 § 2，习题 6）。
\item
  设 G 为一个单的（§ 4，第 4 号，定义 7）带算子交换群。证明 G 的自同态环 A 是一个域。
\item
  设 K 为一个域。证明存在 K 的一个最小子域 F，并且 F 典范同构于域 Q 或域 Z/pZ（其中 p 为某个素数）。在前一种情形（相应地，后一种情形）下，称 K 的特征为 0（相应地，p）。
\item
  设 K 为一个特征 \(\neq2\) 的交换域；设 G 为 K 的加法群的一个子群，使得若 H 表示由 0 与 G 中元素 \(\neq0\) 的逆元组成的集合，则 H 也是 K 的加法群的一个子群。证明存在一个元素 \(a\in K\) 与 K 的一个子域 \(K'\) ，使得 \(G=aK'\) （首先证明：若 x 与 y 是 G 的元素且 \(y\neq0\) ，则 \(x^{2}/y\in G\) ；由此推出：若 x,y,z 是 G 的元素且 \(z\neq0\) ，则 xy/z\ensuremath{\in}G）。
\item
  设 K 为一个特征 \(\neq2\) 的交换域；设 f 为 K 到 K 中的一个映射，使得 \(f(x+y)=f(x)+f(y)\) 对所有 x 与 y 成立，且对所有 \(x\neq0\) 有 \(f(x)f(1/x)=1\) 。证明 f 是 K 到 K 的一个子域上的同构（证明 \(f(x^{2})=(f(x))^{2}\)）。
\item
  设 A 是一个交换环。
\end{enumerate}

\begin{enumerate}
\def\labelenumi{(\alph{enumi})}
\item
  每个素理想都是不可约的（§ 8，习题 11）。
\item
  素理想的集合对于关系 \(\subset\) 与 \(\supset\) 是归纳的。
\item
  设 \(\mathfrak{a}\) 为 \(\mathbf{A}\) 的一个异于 \(\mathbf{A}\) 的理想；令 \(\mathfrak{b}\) 为由这样的 \(x \in \mathbf{A}\) 组成的集合：存在一个整数 \(n \geqslant 0\) 使得 \(x^n \in \mathfrak{a}\) （\(n\) 依赖于 \(x\)）。证明 \(\mathfrak{b}\) 是一个理想，且等于 \(\mathbf{A}\) 的包含 \(\mathfrak{a}\) 的素理想之交。
\end{enumerate}

¶ 8. 环 A 称为布尔环，如果它的每个元素都是幂等的（换句话说，如果 \(x^2 = x\) 对所有 \(x \in \mathbf{A}\) 成立）。

\begin{enumerate}
\def\labelenumi{(\alph{enumi})}
\item
  在集合 \(\mathfrak{P}(\mathbf{E})\) （集合 \(\mathbf{E}\) 的子集组成的集合）中，证明通过令 \(\mathbf{AB} = \mathbf{A} \cap \mathbf{B}\) 与 \(\mathbf{A} + \mathbf{B} = (\mathbf{A} \cap \mathbb{C}\mathbf{B}) \cup (\mathbf{B} \cap \mathbb{C}\mathbf{A})\) 定义了一个布尔环结构。这个环同构于环 \(\mathbf{K}^{\mathrm{E}}\) ——\(\mathbf{E}\) 到模 2 的整数环 \(\mathbf{K} = \mathbf{Z}/(2)\) 中的映射所成的环（对每个子集 \(\mathbf{X} \subset \mathbf{E}\) ，考虑其“特征函数” \(\phi_{\mathbf{X}}\) ，满足 \(\phi_{\mathbf{X}}(x) = 1\) （若 \(x \in \mathbf{X}\)），\(\phi_{\mathbf{X}}(x) = 0\) （若 \(x \notin \mathbf{X}\)））。
\item
  每个布尔环 \(\mathbf{A}\) 都是交换的，且满足 \(x + x = 0\) （对 \(x \in \mathbf{A}\)）（把 \(x + x\) 写成幂等元，然后把 \(x + y\) 写成幂等元）。
\item
  如果布尔环 \(\mathbf{A}\) 不含零因子，则它或者化为 0，或者同构于 \(\mathbf{Z} / (2)\) （若 \(x\) 与 \(y\) 是 \(\mathbf{A}\) 的任意两个元素，证明 \(xy(x + y) = 0\)）。由此推出，在布尔环中每个素理想都是极大理想。
\item
  在布尔环 A 中，每个理想 \(a \neq A\) 都是包含 a 的素理想之交（应用习题 7(c)）。由此推出每个不可约理想都是极大理想（换句话说，在布尔环中，不可约理想、素理想与极大理想这三个概念是一致的）。
\item
  证明每个布尔环都同构于乘积环 \(\mathbf{K}^{\mathbb{E}}\) 的一个子环，其中 \(\mathbf{K} = \mathbf{Z}/(2)\) （利用 \((d)\)）。
\item
  若 \(p_{i}\) ( \(1 \leqslant i \leqslant n\) ) 是布尔环 A 中 n 个互不相同的极大理想，且 \(a = \bigcap_{1 \leqslant i \leqslant n} p_{i}\) ，证明商环 A/a 同构于乘积环 \(K^{n}\) 。由此推出每个有限布尔环都具有形式 \(K^{n}\) 。
\item
  在布尔环 A 中，关系 \(xy = x\) 是一个序关系；若把它记作 \(x \leqslant y\) ，证明 A 在此关系下是一个分配格（集合论，III，§ 1，习题 16），有最小元素 \(\alpha\) ，且对每个有序元素对 \((x, y)\) （满足 \(x \leqslant y\) ），存在一个元素 \(d(x, y)\) ，使得 \(\inf(x, d(x, y)) = \alpha\) , \(\sup(x, d(x, y)) = y\) 。反之，若一个有序集 A 具有这三条性质，证明合成律 \(xy = \inf(x, y)\) 与 \(x + y = d(\sup(x, y), \inf(x, y))\) 在 A 上定义了一个布尔环结构。
\end{enumerate}

\begin{enumerate}
\def\labelenumi{\arabic{enumi}.}
\setcounter{enumi}{8}
\tightlist
\item
  设 \(\mathbf{A}\) 为一个环，使得 \(x^3 = x\) 对所有 \(x \in \mathbf{A}\) 成立。我们要证明 \(\mathbf{A}\) 是交换的。
\end{enumerate}

\begin{enumerate}
\def\labelenumi{(\alph{enumi})}
\item
  证明 \(6\mathbf{A} = \{0\}\) ，且 \(2\mathbf{A}\) 与 \(3\mathbf{A}\) 是双边理想，满足 \(2\mathbf{A} + 3\mathbf{A} = \mathbf{A}\) 且 \(2\mathbf{A} \cap 3\mathbf{A} = \{0\}\) 。由此推出，可以假设或者 \(2\mathbf{A} = \{0\}\) ，或者 \(3\mathbf{A} = \{0\}\) 。
\item
  若 2A = \{0\}，计算 \((1 + x)^{3}\) ，推出 \(x^{2} = x\) 对所有 \(x \in A\) 成立，并利用习题 8 得出结论。
\item
  若 \(3\mathbf{A} = \{0\}\) ，计算 \((x + y)^3\) 与 \((x - y)^3\) ，证明 \(x^2y + xyx + yx^2 = 0\) 并用 \(x\) 左乘；由此推出 \(xy - yx = 0\) 。
\item
  设 A 为一个环，满足 \(3A = \{0\}\) 且 \(x^{3} = x\) 对所有 \(x \in A\) 成立；定义一个集合 I 以及 A 到 \((\mathbf{Z}/3\mathbf{Z})^{\mathrm{I}}\) 中的单射同态（方法与习题 8 相同）。
\end{enumerate}

\begin{enumerate}
\def\labelenumi{\arabic{enumi}.}
\setcounter{enumi}{9}
\tightlist
\item
  设 \(\mathbf{A}\) 是一个交换环， \(\mathfrak{U}\) 是其极大理想的交集。
\end{enumerate}

\begin{enumerate}
\def\labelenumi{(\alph{enumi})}
\item
  证明典范同态 \(\mathbf{A}^* \to (\mathbf{A} / \mathfrak{U})^*\) 是满射的，且其核为 \(1 + \mathfrak{U}\) 。
\item
  假设 A 的极大理想之集 M 是有限的，且 A* 是有限的。由此推出 A 是有限的（把 § 8 第 10 号的命题 8 应用于 A/U）。
\item
  由此推出素数集合是无限的（注意 \(\mathbf{Z}^{*} = \{1, -1\}\)）。
\item
  假设 A* 是有限的，且对每个极大理想 m，A/m 都是有限的。能否由此推出 A 是有限的？A 是可数的？
\end{enumerate}

\begin{enumerate}
\def\labelenumi{\arabic{enumi}.}
\setcounter{enumi}{10}
\tightlist
\item
  设 P 为素数集合，A 为域 \(\mathbf{Z} / p\mathbf{Z}\) , \(p \in \mathbb{P}\) 的乘积环。设 \(a\) 为 A 的由元素 \((a_p)_{p \in \mathbb{P}}\) （满足 \(a_p \neq 0\) 仅对有限多个指标 \(p\) 成立）组成的子集。
\end{enumerate}

\begin{enumerate}
\def\labelenumi{(\alph{enumi})}
\item
  a 是 A 的一个理想。
\item
  设 \(\mathbf{B} = \mathbf{A} / a\) 。对每个整数 \(n > 0\) 与每个元素 \(b \neq 0\) （在 \(\mathbf{B}\) 中），存在且仅存在一个元素 \(b'\) （属于 \(\mathbf{B}\)），使得 \(nb' = b\) （注意，若 \(p \in \mathbb{P}\) 不整除 \(n\) ，则用 \(n\) 乘在 \(\mathbf{Z} / p\mathbf{Z}\) 中是双射）。由此推出 \(\mathbf{B}\) 包含一个同构于 \(\mathbf{Q}\) 的子域。
\end{enumerate}

\begin{enumerate}
\def\labelenumi{\arabic{enumi}.}
\setcounter{enumi}{11}
\item
  设 A 为一个交换环，且 \(a, b \in A\) 。证明 ab 在 \(\mathrm{A}/(a - a^{2}b)\) 中的典范像是幂等的。给出一个使该幂等元异于 0 与 1 的例子。
\item
  确定环 Z 与环 \(Z \times Z\) 的自同态。更一般地，若 I 与 J 是有限集，确定环 \(Z^{1}\) 到环 \(Z^{J}\) 中的同态。
\item
  设 A 与 B 是两个环，f 与 g 是 A 到 B 中的两个同态。使 \(x \in A\) 满足 \(f(x) = g(x)\) 者组成的集合是 A 的子环、A 的理想吗？
\end{enumerate}

¶ 15. 设 A 是一个无零因子的非交换环。称 A 容许一个左分式域，如果它同构于某个域 K 的子环 B，使得 K 的每个元素都具有形式 \(x^{-1}y\) ，其中 \(x \in \mathbf{B}, y \in \mathbf{B}\) 。

\begin{enumerate}
\def\labelenumi{(\alph{enumi})}
\tightlist
\item
  设 A' 为 A 中元素 \(\neq0\) 组成的集合。为使 A 容许一个左分式域，下列条件必须满足：
\end{enumerate}

\begin{enumerate}
\def\labelenumi{(\Alph{enumi})}
\setcounter{enumi}{6}
\tightlist
\item
  对所有 \(x \in \mathbf{A}\) , \(x' \in \mathbf{A}'\) ，存在 \(u \in \mathbf{A}'\) 与 \(v \in \mathbf{A}\) ，使得 \(ux = vx'\) 。
\end{enumerate}

\begin{enumerate}
\def\labelenumi{(\alph{enumi})}
\setcounter{enumi}{1}
\tightlist
\item
  反之，假设条件 (G) 成立。证明在集合 A × A' 中，\((x, x')\) 与 \((y, y')\) 之间的关系 R——即“对每个由非零元素组成的有序对 \((u, v)\) ，只要 \(ux' = vy'\) ，就有 \(ux = vy\) ”——是一个等价关系。
\end{enumerate}

设 \((x, x')\) , \((y, y')\) 为 \(A \times A'\) 的两个元素，\(\xi\) 与 \(\eta\) 为它们各自的类（mod. R）。对每个有序对 \((u, u') \in A \times A'\) （满足 \(u'x = uy'\) ），证明 \((uy, u'y')\) 的类（mod. R）仅依赖于类 \(\xi\) 与 \(\eta\) ；若把它记作 \(\xi\eta\) ，就在集合 \(K = (A \times A')/R\) 上定义了一个合成律。若 \(K'\) 为 K 中异于诸元素 \((0, x')\) （属于 \(A \times A'\)）的类 0 的元素组成的集合，则 \(K'\) 连同由上述法则诱导的法则构成一个群。

对每个元素 \(x \in \mathbf{A}\) ，诸元素 \((x'x, x')\) （其中 \(x'\) 遍历 \(\mathbf{A}'\) ）定义了 \(\mathbf{A}\) （仅带乘法）到 \(\mathbf{K}\) 的一个子环上的同构。把 \(\mathbf{A}\) 与它在此同构下的像等同以后，有序对 \((x, x') \in \mathbf{A} \times \mathbf{A}'\) 的类（mod. R）就与元素 \(x'^{-1}x\) 等同。

既然如此，若 \(\xi = x'^{-1}x\) 是 K 的一个元素，且 1 表示 \(\mathbf{K}'\) 的单位元素，则 \(\xi + 1\) 表示元素 \(x'^{-1}(x + x')\) ，它不依赖于 \(\xi\) 写成形式 \(x'^{-1}x\) 的方式。然后我们写 \(\xi + 0 = \xi\) ，且对 \(\eta \neq 0\) ，写 \(\xi + \eta = \eta(\eta^{-1}\xi + 1)\) 。证明这样在 K 上定义的加法与乘法在该集合上确定了一个域结构，它扩充了 A 上的环结构；换言之，条件 (G) 使 A 容许一个左分式域是充分的。

\begin{enumerate}
\def\labelenumi{\arabic{enumi}.}
\setcounter{enumi}{15}
\item
  设 A 是一个环，其中每个非零元素都是可消去的。为使 A 容许一个左分式域（习题 15），必要且充分的是：在 A 中，两个异于 0 的左理想之交永不缩减为 0。
\item
  设 \((\mathbf{K}_i)_{i\in I}\) 是一个域族。对每个子集 \(J\) ⊂ \(I\) ，设 \(e_J\) 表示乘积 \(\mathbf{K} = \prod_{i\in I}\mathbf{K}_i\) 中这样的元素：它的第 \(i\) 个分量当 \(i\in J\) 时等于 0，而当 \(i\in I - J\) 时等于 1。
\end{enumerate}

\begin{enumerate}
\def\labelenumi{(\alph{enumi})}
\item
  设 \(x = (x_{i})\) 是 \(\mathbf{K}\) 的一个元素，\(J(x)\) 为元素 \(i \in I\) 中使 \(x_{i} = 0\) 者组成的集合。证明存在一个可逆元素 \(u\) （在 \(\mathbf{K}\) 中），使得 \(x = ue_{J(x)} = e_{J(x)}u\) 。
\item
  设 \(\mathfrak{a}\) 是 \(\mathbf{K}\) 的一个异于 \(\mathbf{K}\) 的左（相应地，右）理想。设 \(\mathfrak{F}_{\mathfrak{a}}\) 为子集 \(\mathbf{J}\) （属于 \(\mathbf{I}\) 且满足 \(e_{J} \in \mathfrak{a}\)）组成的集合。证明 \(\mathfrak{F}_{\mathfrak{a}}\) 具有下列性质：
\end{enumerate}

\((b_{1})\) 若 \(\mathbf{J} \in \mathfrak{F}_{\mathfrak{a}}\) 且 \(\mathbf{J}' \supset \mathbf{J}\) ，则 \(\mathbf{J}' \in \mathfrak{F}_{\mathfrak{a}}\) 。

\((b_{2})\) 若 \(\mathrm{J}_1, \mathrm{J}_2 \in \mathfrak{F}_{\mathfrak{a}}\) ，则 \(\mathrm{J}_1 \cap \mathrm{J}_2 \in \mathfrak{F}_{\mathfrak{a}}\) 。

\[
\left(b _ {3}\right) \varnothing \notin \mathfrak {F} _ {\mathfrak {a}}.
\]

证明 \(x \in \mathfrak{a}\) 等价于 \(J(x) \in \mathfrak{F}_{\mathfrak{a}}\) （利用 (a)）。由此推出 \(\mathfrak{a}\) 是一个双边理想。

\begin{enumerate}
\def\labelenumi{(\alph{enumi})}
\setcounter{enumi}{2}
\item
  I 的一个子集族称为滤子，如果它满足上述性质 \((b_{1})\) , \((b_{2})\) , \((b_{3})\) （参见《一般拓扑学》I，§6）。证明映射 \(\mathfrak{a} \mapsto \mathfrak{F}_{\mathfrak{a}}\) 是 K 的异于 K 的理想之集到 I 的滤子之集上的一个严格递增双射。*证明 a 是极大的当且仅当 \(\mathfrak{F}_{\mathfrak{a}}\) 是一个超滤子。*
\item
  *设 \(\mathfrak{F}\) 是素数集合 \(P\) 上的一个非平凡超滤子，\(a\) 为环 \(\mathbf{K} = \prod_{p \in P} \mathbf{Z}/p\mathbf{Z}\) 的相应理想，并设 \(k = K/a\) 。证明 \(k\) 是一个特征为 \(0\) 的域。*
\end{enumerate}

¶ 18. (a) 设 K 是一个域，a、b 是 K 中两个不可交换的元素。证明

\[
\begin{array}{l} (1) a = (b - (a - 1) ^ {- 1} b (a - 1)) (a ^ {- 1} b a - (a - 1) ^ {- 1} b (a - 1)) ^ {- 1} \\ (2) a = (1 - (a - 1) ^ {- 1} b ^ {- 1} (a - 1) b) (a ^ {- 1} b ^ {- 1} a b - (a - 1) ^ {- 1} b ^ {- 1} (a - 1) b) ^ {- 1}. \end{array}
\]

\begin{enumerate}
\def\labelenumi{(\alph{enumi})}
\setcounter{enumi}{1}
\item
  设 K 是一个非交换域，x 是不属于 K 的中心的一个元素。证明 K 由 x 的共轭元素 \(axa^{-1}\) 组成的集合生成。（设 \(K_{1}\) 为 K 的由 x 的共轭元素集合生成的子域。由 (a) 推出，若 \(K_{1} \neq K\) 且 \(a \in K_{1}\) 而 \(b \notin K_{1}\) ，则必有 ab = ba。通过在 \(K_{1}\) 中考虑两个不可交换的元素 a、\(a'\) 以及一个元素 \(b \notin K_{1}\) ，并注意到 \(ba \notin K_{1}\) ，由此得出矛盾。）
\item
  由 (b) 推出，若 Z 是 K 的中心且 H 是 K 的一个子域，使得 \(Z \subset H \subset K\) 且 \(aHa^{-1} = H\) 对 K 中所有 \(a \neq 0\) 成立，则 H = Z 或 H = K（嘉当–布劳尔–华定理）。
\item
  由 (a) 推出，在一个非交换域 K 中，换位子 \(aba^{-1}b^{-1}\) （元素 \(\neq0\) 的换位子）所组成的集合生成域 K。
\end{enumerate}

\begin{enumerate}
\def\labelenumi{\arabic{enumi}.}
\setcounter{enumi}{18}
\tightlist
\item
  设 A 是一个非零伪环，使得对 A 中所有 \(a \neq 0\) 及所有 \(b \in \mathbf{A}\) ，方程 \(ax + ya = b\) 都有由 A 的元素组成的有序对 \((x, y)\) 作为解；换言之，\(a\mathbf{A} + \mathbf{A}a = \mathbf{A}\) 对 A 中的 \(a \neq 0\) 成立。
\end{enumerate}

\begin{enumerate}
\def\labelenumi{(\alph{enumi})}
\tightlist
\item
  证明 A 是 A 中唯一的非零双边理想。
\end{enumerate}

\begin{enumerate}
\def\labelenumi{(\alph{enumi})}
\setcounter{enumi}{1}
\item
  证明关系 \(a \neq 0\) 蕴含 \(a^2 \neq 0\) （注意，若 \(a^2 = 0\) ，则将得出 \(aAa = \{0\}\) ）。
\item
  证明：若 \(ab = 0\) 且 \(a \neq 0\) （相应地，\(b \neq 0\) ），则 \(b = 0\) （相应地，\(a = 0\) ）。（注意 \((ba)^2 = 0\) ，并由此推出 \(bAb = \{0\}\) 。）
\end{enumerate}

\begin{enumerate}
\def\labelenumi{\arabic{enumi}.}
\setcounter{enumi}{19}
\item
  设 A 是一个非零伪环，使得存在 \(a \in A\) ，对每个 \(b \in A\) ，方程 ax = b、xa = b 中有一个具有解 \(x \in A\) 。证明：若 A 除 0 外没有零因子，则 A 具有单位元素。
\item
  设 A 是一个非零伪环，其中对所有 \(a \neq 0\) 以及所有 \(b \in A\) ，方程 ax = b、xa = b 中有一个具有解 \(x \in A\) 。证明 A 是一个域（利用习题 19 和 20）。
\end{enumerate}

\BourbakiExercisesSection

\begin{enumerate}
\def\labelenumi{\arabic{enumi}.}
\tightlist
\item
  设 X 为集合的正向系统 \((\mathbf{X}_{i}, f_{ji})\) 关于右有向集 I 的正向极限。设 \(F_{i}\) （相应地，F）为在 \(X_{i}\) （相应地，X）上构造的自由群；每个 \(f_{ji}\) 可扩张为一个同态 \(\phi_{ji}: F_{i} \to F_{j}\) 。
\end{enumerate}

\begin{enumerate}
\def\labelenumi{(\alph{enumi})}
\item
  证明 \((\mathbf{F}_i,\phi_{ji})\) 是一个群的正向系统。
\item
  设 \(\varepsilon_{i}\) 为复合映射 \(\mathbf{X}_{i} \to \mathbf{F}_{i} \to \lim \mathbf{F}_{i}\) 。证明 \(\varepsilon_{j} \circ \phi_{ji} = \varepsilon_{i}\) 当 \(j \geqslant i\) 时成立。设 \(\varepsilon\) 为 \(\mathbf{X}\) 到 \(\varinjlim \mathbf{F}_{i}\) 中的、由诸 \(\varepsilon_{i}\) 所定义的映射。证明 \(\varepsilon\) 可扩张为 \(\mathbf{F}\) 到 \(\varinjlim \mathbf{F}_{i}\) 上的一个同构（因此 \(\mathbf{F}(\varinjlim \mathbf{X}_{i})\) 可与 \(\lim \mathbf{F}(\mathbf{X}_{i})\)
\item
  等同）。对自由广群、自由幺半群、自由交换群、自由交换幺半群，陈述并证明类似的结果。
\end{enumerate}

\begin{enumerate}
\def\labelenumi{\arabic{enumi}.}
\setcounter{enumi}{1}
\tightlist
\item
  证明单群的正向极限或者是单群，或者是只有一个元素的群。
\end{enumerate}

\BourbakiStarSection{历史注记}

（方括号中的数字指本注末尾的参考文献。）

数学中几乎没有比合成律更原始的概念了：它似乎与自然数和可测量量的最初计算基础密不可分。流传至今的关于埃及人和巴比伦人数学的最古老文献表明，他们早已拥有了一套关于自然数 \textgreater0 、有理数 \textgreater0 、长度与面积的完整的计算规则体系；尽管保存下来的文本仅涉及给定数量具有明确数值的问题，† 但它们使我们毫不怀疑他们对所用规则所赋予的普遍性，并展示了在处理一次和二次方程方面相当卓越的技术能力（{[}1{]}，第 179 页及以下）。另一方面，完全没有丝毫关注去证明所使用的规则，甚至对所出现的运算也没有精确的定义：后者与前者一样，都是以纯粹经验的方式产生的。

然而，正是这样一种关注，在古典时代的希腊人的著作中已经非常清晰地显现出来；诚然，自然数理论并未找到公理化的处理（这样的公理化直到十九世纪末才出现；参见《集合论》第 IV 章的历史注记）；但在欧几里得的《几何原本》中，有许多段落对计算规则给出了形式化的证明，这些规则与整数计算的规则一样，在直觉上是“显然”的（例如，两个有理数乘积的交换性）。这类性质最引人注目的证明，是关于量理论的证明，这是希腊数学最具原创性的创造（众所周知，它相当于我们的实数理论 \textgreater0 ；参见《一般拓扑学》第 IV 章的历史注记）；在这里，欧几里得除其他事项外，考虑了两种量之比（ratio）的乘积，并表明该乘积与这些比的表示形式无关（这是 § 1 第 6 号中按等价关系对合成律取“商”的第一个例子），并且这些乘积满足交换律（{[}2{]}，第五卷，命题 22—23）。†

然而，不可隐瞒的是，这种向严格性的进步在欧几里得那里伴随着停滞，甚至在代数计算技术方面在某些方面还有所倒退。几何的压倒性优势（显然是在其光照下构思了量论的）使得代数符号的任何独立发展都陷入瘫痪：进入计算的元素必须始终在几何上被“表示”；此外，出现的两种合成律并非定义在同一集合上（比值的加法并未以一般方式定义，而两条长度的乘积不是长度而是面积）；其结果是缺乏灵活性，使得操纵次数大于二的代数关系几乎不可行。

只有当古典希腊数学衰落时，丢番图才回归到“计算师”或专业计算者的传统，这些人一直沿用从埃及人和巴比伦人那里继承来的规则：他不再用几何表示来束缚自己所考虑的“数”，自然便发展出抽象代数计算的规则；例如，他给出的规则（用现代语言来说）对 m 和 n 的小值（正或负）等价于公式 \(x^{m+n} = x^{m}x^{n}\) （{[}3{]}，第 I 卷，第 8—13 页）；稍后（第 12—13 页），陈述了“符号法则”，即负数计算的初步内容‡；最后，丢番图首次使用一个字母来表示方程中的未知数。相比之下，他似乎很少关心为他解决问题所应用的方法赋予普遍意义；至于由欧几里得开创的合成律的公理化观念，这对丢番图及其直接后继者的思想而言似乎是陌生的；这种观念直到 19 世纪初才在代数学中重新出现。

在随后的几个世纪中，首先需要发展一套足以表达抽象规律的代数符号体系，另一方面则需要通过观察足够多样的特例，将“数”的概念加以拓展，从而形成一般性的观念。为此，希腊人创立的量之比的公理化理论并不足够，因为它只不过使实数 \textgreater0 的直观概念以及这些数上的运算变得精确而已，而这些运算巴比伦人早已以较为混乱的形式知晓；如今所考虑的“数”将是希腊人从未有过的概念，并且起初也没有可感知的“表示”：一方面是零和负数，它们出现在中世纪盛期的印度数学中；另一方面是虚数，这是 16 世纪意大利代数学家的创造。

除了零——最初是作为分隔符号引入，后来才被视为数（见《集合论》第 III 章的历史注记）——这些扩充的共同特点是（至少起初）它们纯粹是“形式的”。这里必须理解为，新的“数”首先表现为对某些情形进行运算的结果，而在这些情形中，若严格遵循其定义，这些运算本无意义（例如，两个自然数之差 \(a - b\) 当 \(a < b\) 时的情形）：由此产生了赋予它们的名称：“虚假的”、“虚构的”、“荒谬的”、“不可能的”、“想象的”数等等。对于古典时期的希腊人而言，他们尤其热爱清晰的思想，这样的扩展是不可想象的；它们只能出现在那些比希腊人更倾向于对其方法的力量表现出某种神秘信仰（如 18 世纪所说的“分析的普遍性”），并任由自己随计算的机械过程前行而不去探究每一步是否有充分依据的计算者手中；这种信心事后通常被证明是合理的，因为把这些严格说来仅对已知数有效的计算规则应用于这些新的数学实体，导出了精确的结果。这就解释了为何这些数字概念的推广——起初仅作为一系列运算中的中介出现，其起点和终点都是真正的“数字”——逐渐被独立看待（不依赖于具体计算的应用）；一旦迈出这一步，人们便寻求或多或少可感知的解释，从而赋予新对象在数学中存在的权利。†

在这个问题上，印度人已经意识到在某些情况下负数应被赋予的解释（例如商业问题中的债务）。在随后的几个世纪里，随着希腊和印度数学的方法与成果（经由阿拉伯人）传播到西方，这些数的运算变得更加熟悉，并获得了其他具有几何或运动学性质的“表示”。而代数符号体系的逐步改进，则是中世纪末期代数学中唯一显著的进步。

16 世纪初，由于意大利学派数学家发现了三次方程及随后四次方程的“根式”解法（我们将在《代数》第 V 章的历史注记中更详细地讨论这一点），代数学获得了新的推动力；正是在此时，尽管他们心存厌恶，却不得不将虚数引入计算之中；此外，对这些“不可能”数的运算信心逐渐建立起来，正如对负数的运算一样，尽管在两个多世纪里，人们从未构想出它们的任何“表示”。

另一方面，代数符号体系由韦达和笛卡尔决定性地完善；自笛卡尔之后，代数符号大体上就是我们今天所使用的形式。

从 17 世纪中叶到 18 世纪末，微积分的创立所开辟的广阔视野似乎导致了代数学总体上的一定忽视，尤其是对合成律以及实数与复数性质的数学反思。† 因此，自 17 世纪末在力学中已广为人知的力的合成与速度的合成，在代数学中并未产生反响，尽管它们已经蕴含了向量演算的萌芽。事实上，必须等到大约 1800 年前后，导致复数几何表示的思想运动（见《一般拓扑学》第 VIII 章的历史注记），才看到向量加法在纯数学中得到应用。‡

大约在同一时期，代数学中首次将合成律的概念沿两个不同方向推广到与“数”（迄今为止该词最广义的含义）仅有遥远类比关系的元素上。这两个推广中的第一个归功于 C. F. 高斯，起因于他对二次型 \(ax^{2} + bxy + cy^{2}\) （具有整数系数）的算术研究。拉格朗日曾在具有相同判别式的形式集合上定义了一个等价关系†，并且另一方面证明了一个恒等式，该恒等式为该集合提供了一个交换合成律（并非处处有定义）；基于这些结果，高斯证明了该合成律与上述等价关系的一种细化形式相容（在 § 1 第 6 号的意义下）（{[}5{]}，第 I 卷，第 272 页）：“这表明”，他随后说道，“必须理解两个或若干类的复合类是什么意思”。接着他继续研究“商”律，实质上确立了它（用现代语言说）是一个交换群律，并且所用的论证其普遍性在很大程度上远远超出了高斯所考虑的特殊情形（例如，他证明逆元唯一性的论证与我们针对 § 2 第 3 号中任意合成律所给出的论证完全相同（同上，第 273 页））。他并未止步于此：稍后回到这个问题时，他认识到类的复合与模素数整数的乘法之间的类比‡（同上，第 371 页），但也证明了具有给定判别式的二次形式的类群并不总是循环群；他在这方面给出的指示清楚地表明，他至少在这个特殊情形下已经认识到了有限阿贝尔群的一般结构，我们将在第七章中研究这一结构（{[}5{]}，第 I 卷，第 374 页及第 II 卷，第 266 页）。

我们想谈的另一系列研究也以群的概念及其最终引入数学而告终：这就是“代换理论”，即拉格朗日、范德蒙德和高斯关于代数方程求解思想的发展。我们在此不给出这一主题的详细历史（参见《代数》第 V 章的历史注记）；我们必须提及鲁菲尼以及随后柯西（{[}6{]}，（2），第 I 卷，第 64 页）对有限集合的两个置换的“乘积”的定义§，以及关于有限置换群的初步概念：传递性、本原性、单位元、可交换元素等。但这些初步结果总体上仍然相当肤浅，埃瓦里斯特·伽罗瓦必须被视为该理论的真正开创者：在他那篇令人难忘的著作 {[}8{]} 中，他将代数方程的研究归结为与之相关的置换群的研究，从而大大深化了对后者的研究，无论是在群的一般性质方面（伽罗瓦是第一个定义正规子群的概念并认识到其重要性的人），还是在具有特殊性质的群的确定方面（他在这方面获得的结果至今仍被认为是该理论中最精妙的结果之一）。“群的线性表示”的最初想法也追溯到伽罗瓦†，这一事实清楚地表明他拥有两个群结构的同构概念，而不依赖于它们的“实现”。

然而，尽管高斯和伽罗瓦的巧妙方法无可争辩地引导他们形成了对合成律概念的非常宽广的理解，但他们并未特别在这方面发展他们的思想，他们的工作对抽象代数的演进也没有产生直接的影响。‡ 朝向抽象的最清晰进展是在第三个方向上取得的：在反思虚数的性质之后（虚数的几何表示在 19 世纪初引发了大量的工作），英国学派的代数学家们率先在 1830 年至 1850 年间分离出合成律的抽象概念，并立即通过将这一概念应用于大量新的数学实体而拓宽了代数的领域：布尔那里的逻辑代数（参见《集合论》第 IV 章的历史注记），哈密顿那里的向量和四元数，凯莱那里的一般超复数系统、矩阵与非结合律（{[}10{]}，第 I 卷，第 127 页和第 301 页，第 II 卷，第 185 页和第 475 页）。在欧洲大陆也独立地进行了平行的演进，特别是在向量演算（莫比乌斯、贝拉维蒂斯）、线性代数和超复数系统（格拉斯曼）方面，我们将在《代数》第 III 章的历史注记§中更详细地讨论这些内容。

从所有这些在 19 世纪上半叶使代数焕发活力的原创而富有成果的思想的激荡中，代数彻底地焕然一新。此前，方法和结果都围绕着代数方程（或数论中的丢番图方程）求解这一中心问题而展开：“代数”，塞雷特在其《高等代数教程》的引言中说道 {[}12{]} ，“严格说来，是方程的分析”。1850 年之后，尽管代数方面的论著在一段时间内仍然给予方程理论以优先地位，但新的研究工作不再受制于对数值方程求解的直接应用的关注，而是越来越多地朝向今天被认为是代数的本质问题的方向，即为研究而研究代数结构。

这些工作相当整齐地分为三大主流，它们分别延续了我们上面分析过的三个思想运动，并且在19世纪最后几年之前一直平行发展，彼此之间没有显著的影响。†

这些工作主流中的第一个源于 19 世纪德国学派（狄利克雷、库默尔、克罗内克、戴德金、希尔伯特）在建立代数数论方面的工作，其发端于高斯的研究，高斯开创了这类数的第一种研究，即数 \(a + bi\) （ \(a\) 和 \(b\) 为有理数）的研究。我们在此不追踪这一理论的发展（参见《交换代数》第七章的历史注记）：就我们的目的而言，只需提及由此产生的抽象代数概念。从高斯的最初后继者起，域（代数数域）的概念就是这一主题所有著作的基础（正如它对阿贝尔和伽罗瓦关于代数方程的研究一样）；当戴德金和韦伯 {[}13{]} 以代数数理论为模型建立了单变量代数函数理论时，它的应用范围得到了扩大。戴德金 {[}14{]} 还提出了理想的概念，这一概念为元素集合之间的合成律提供了新的范例；他和克罗内克共同推动了交换群和模在代数域理论中日益重要的作用；我们将在《代数》第三、五、七章的历史注记中再次讨论这一点。

关于线性代数与超复数系统的发展，我们还将参考《代数》第 III 章和第 VIII 章的历史注记；这一发展在 19 世纪末至 20 世纪初并未引入任何新的代数概念，它在英国（西尔维斯特、W. 克利福德）和美国（B. 皮尔斯和 C. S. 皮尔斯、迪克森、韦德伯恩）沿着哈密顿和凯莱所开辟的道路进行，同时在德国（魏尔斯特拉斯、戴德金、弗罗贝尼乌斯、莫利恩）和法国（拉盖尔、E. 嘉当）独立于盎格鲁-撒克逊学派，并采用了截然不同的方法。

就群论而言，它最初主要是作为有限置换群的理论发展起来的，这发生在伽罗瓦的著作出版之后，并通过塞雷特的著作 {[}12{]} 以及尤其是 C. 若尔当的伟大著作《置换论》 {[}15{]} 的传播而实现。若尔当在那里极大地改进了其前辈关于置换群特殊性质（传递性、本原性等）的研究成果，获得了大部分至今未被超越的结果；他还对一类非常重要的特殊群——线性群及其子群——进行了同样深入的研究；此外，正是他引入了从一个群到另一个群的同态这一基本概念，并且（稍后）引入了商群的概念，还证明了被称为“若尔当–赫尔德定理”的定理的一部分。† 最后，若尔当是第一个研究无限群的人 {[}16{]} ，这些无限群后来由 S. 李一方面以及 F. 克莱因和 H. 庞加莱另一方面在几年后沿着两个不同方向进行了相当大的发展。

与此同时，群中本质的东西，即其合成律而非构成群的对象的性质，逐渐被认识到（例如参见 {[}10{]}，第 II 卷，第 123 页和第 131 页以及 {[}14{]}，第 III 卷，第 439 页）。然而，即使是对有限抽象群的研究工作，在很长一段时间内也被视为置换群的研究，直到大约 1880 年，有限群的独立理论才开始有意识地发展起来。我们无法进一步追溯这一理论的历史，本专著中对此仅作了非常肤浅的涉及；我们只提及两个至今在研究有限群中仍最常用的工具，它们都追溯到 19 世纪：关于 p 群的西罗定理‡（其年代为 1872 年 {[}17{]} ）以及该世纪末由弗罗贝尼乌斯创立的特征标理论 {[}19{]} 。我们建议希望深入研究有限群理论及其引发的许多难题的读者，参阅伯恩赛德 {[}20{]} 、施派泽 {[}23{]} 、扎森豪斯 {[}24{]} 和戈伦斯坦 {[}27{]} 的专著。

大约在 1880 年，群表现的系统研究也开始了；此前仅遇到过特定群的表现，例如交错群 \(\mathfrak{A}_5\) 在哈密顿的一篇著作中 {[}9 bis{]} 的表现，或线性微分方程和黎曼曲面的单值群（施瓦茨、克莱因、施莱夫利）的表现。W. 迪克是第一个 {[}18{]} 定义（但未给它命名）由有限个生成元生成的自由群的人，但他对自由群本身的兴趣不如把它作为“万有”对象那样大，后者使得可以精确定义“由生成元和关系”给出的群。迪克发展了凯莱首先提出、同时又明显受到上述黎曼学派著作影响的思想，描述了具有给定表现的群的一种解释，其中每个生成元由关于相切或相交圆的两个反演之积来表示（另见伯恩赛德 {[}20{]} ，第十八章）。稍后，在庞加莱及其后继者发展了自守函数理论之后，庞加莱又引入了代数拓扑的工具，对基本群的初步研究便与群表现的研究（德恩、蒂策）并驾齐驱，两种理论相互支持。此外，正是拓扑学家 J. 尼尔森在 1924 年对自由群性质的第一个深入研究中引入了“自由群”这一术语 {[}21{]} ；紧接着，E. 阿廷（始终参照拓扑学中的问题）引入了群的自由积的概念，而 O. 施赖尔则更一般地定义了带融合子群的自由积 {[}22{]} （参见 I，§ 7，第 3 号，习题 20）。在此我们同样无法详述这一方向后续发展的历史，请读者参阅著作 {[}26{]} 以获取更多细节。

自 19 世纪末以来，群的概念（以及与之密切相关的“不变量”概念）在分析学、几何学、力学和理论物理学中所取得的非凡成功，已无需再多言。这一概念以及与之相关的代数概念（带算子的群、环、理想、模）同样侵入到此前看似毫无关联的代数学的各个部分，这构成了我们在此追溯的、并以我们上文所追踪的三个分支的综合而告终的最后一个发展时期的显著标志。这一统一主要是 1900 年至 1930 年间德国学派的工作成果：由戴德金和希尔伯特在 19 世纪最后几年开创的代数学公理化工作，由 E. 施泰尼茨有力地继续推进，随后自 1920 年起，在 E. 阿廷、E. 诺特及其学派代数学家（哈塞、克鲁尔、O. 施赖尔、范德瓦尔登）的推动下得到大力发展。范德瓦尔登的专著 {[}25{]} 于 1930 年出版，首次将这些工作汇集于一部统一的论述之中，为二十多年间代数学领域的许多研究工作开辟了道路并起到了指导作用。

\BourbakiStarSection{参考文献}

\begin{enumerate}
\def\labelenumi{\arabic{enumi}.}
\item
  O. 诺伊格鲍尔，《古代数学史讲义》，第 I 卷：《前希腊数学》，柏林（施普林格），1934 年。
\item
  《欧几里得原本》，5卷本，J. L. 海贝格编，莱比锡（托伊布纳），1883—88年。
\end{enumerate}

2之二. T. L. 希思，《欧几里得〈原本〉十三卷》……，3卷本，剑桥，1908年。

\begin{enumerate}
\def\labelenumi{\arabic{enumi}.}
\setcounter{enumi}{2}
\tightlist
\item
  《亚历山大的丢番图全集》 \ldots，2卷本，P. 坦纳里编，莱比锡（托伊布纳），1893—95年。
\end{enumerate}

3之二. 亚历山大的丢番图，P. 韦尔·埃克译，布鲁日（德克莱-德布鲁韦），1926年。

\begin{enumerate}
\def\labelenumi{\arabic{enumi}.}
\setcounter{enumi}{3}
\item
  G. W. 莱布尼茨，《数学著作》，C. I. 格哈特编，第 V 卷，哈雷（施密特），1858 年。
\item
  C. F. 高斯，《全集》，第 I 卷（哥廷根，1870 年）、第 II 卷（同上，1863 年）和第 VIII 卷（同上，1900 年）。
\item
  A. L. 柯西，《全集》（第二辑），第 I 卷，巴黎（戈蒂埃-维拉尔），1905 年。
\item
  N. H. 阿贝尔，《著作集》，2 卷本，西罗和利编，克里斯蒂安尼亚，1881 年。
\item
  E. 伽罗瓦，《数学著作与论文》，R. 布尔涅和J. Y. 阿兹拉编，巴黎（戈蒂埃-维拉尔），1962年。
\item
  W. R. 哈密顿，《四元数讲义》，都柏林，1853年。
\end{enumerate}

9之二. W. R. 哈密顿，《关于单位根的一个新体系的备忘录》，《哲学杂志》（第四辑），12（1856 年），第 446 页。

\begin{enumerate}
\def\labelenumi{\arabic{enumi}.}
\setcounter{enumi}{9}
\item
  A. 凯莱，《数学论文集》，第 I 卷和第 II 卷，剑桥（大学出版社），1889 年。
\item
  H. 汉克尔，《复数及其函数讲义》，第一部分：《复数系理论》，莱比锡（福斯），1867年。
\item
  J. A. 塞雷特，《高等代数教程》，第 3 版，巴黎（戈蒂埃-维拉尔），1866 年。
\item
  R. 戴德金和H. 韦伯，《单变量代数函数理论》，《克雷尔杂志》，92（1882年），第181—290页。
\item
  R. 戴德金，《数学全集》，3卷本，不伦瑞克（菲韦格），1932年。
\item
  C. 若尔当，《置换与代数方程论》，巴黎（戈蒂埃-维拉尔），1870年（重印，巴黎（A. 布朗夏尔），1957年）。
\item
  C. 若尔当，《关于运动群的论文》，《数学纪事》（第二辑），2（1868 年），第 167—215 页及第 322—345 页（=《全集》，第 IV 卷，第 231—302 页，巴黎（戈蒂埃-维拉尔），1964 年）。
\item
  L. 西罗，《关于置换群的定理》，《数学纪事》，5（1872 年），第 584—594 页。
\item
  W. 迪克，《群论研究》，《数学纪事》，20（1882 年），第 1—44 页。
\item
  G. 弗罗贝尼乌斯，《关于群的特征标》，《柏林会议报告》，1896 年，第 985—1021 页（=《论文集》，J. P. 塞尔编，柏林-海德堡-纽约（施普林格），第 III 卷（1968 年），第 1—37 页）。
\item
  W. 伯恩赛德，《有限阶群论》，第2版，剑桥，1911年。
\item
  J. 尼尔森，《自由群的自同构群》，《数学纪事》，91（1924 年），第 169—209 页。
\item
  O. 施赖尔，《自由群的子群》，《汉堡大学数学讨论会报告》，5（1927 年），第 161—185 页。
\end{enumerate}

\begin{enumerate}
\def\labelenumi{\arabic{enumi}.}
\setcounter{enumi}{22}
\item
  A. 施派泽，《有限阶群论》，第3版，柏林（施普林格），1937年。
\item
  H. 扎森豪斯，《群论教科书》，第 I 卷，莱比锡-柏林（托伊布纳），1937 年。
\item
  B. L. 范德瓦尔登，《近世代数》，第2版，第I卷，柏林（施普林格），1937年；第II卷（同上），1940年。
\item
  W. 马格努斯、A. 卡拉斯和 D. 索利塔，《组合群论》，纽约（Interscience），1966 年。
\item
  D. 戈伦斯坦，《有限群》，纽约（哈珀与罗），1968年。
\end{enumerate}

\chapter{线性代数}

本章主要致力于研究一类特殊的带算子的交换群（I，§ 4，第 2 号），称为模。§§ 1 和 2 中为模所陈述的某些性质可推广到所有带算子的交换群；这些性质将在出现时予以指出。此外，将在第三章 § 2 第 6 号中看到，对带算子的交换群的研究总是等价于对与该带算子的群适当关联的一个模的研究。

\section{模}
